\section{5. Experiments — CLEAN REVISED 2026-04-14}
\subsection{5.1 Protocol and reproducibility}

\textbf{Models (FC-native Instruct primary roster).} Tool-selection evaluation is meaningful only on models trained to emit structured function calls. All primary cells use FC-capable Instruct variants with \texttt{tools} support in their chat template:


\begin{table}[t]
\centering
\begin{tabularx}{\linewidth}{X | X | X | X | X | X}
\toprule
\textbf{P1 primary} & \texttt{Qwen/Qwen2.5-7B-Instruct} & ✓ & C & 4 & Main reference; scaling pivot \\
\midrule
P1 primary & \texttt{NousResearch/Meta-Llama-3.1-8B-Instruct} (un-gated mirror) & ✓ & A & 8 & Cross-family (Mode A ✓) \\
P1 primary & \texttt{mistralai/Mistral-7B-Instruct-v0.3} & ✓ & A & 8 & 86/14 counterexample + H2 \\
P1 stretch & \texttt{google/gemma-3-27b-it} (pending gated approval) & ✓ & — & — & \textbf{Netsru deployment model} — direct production alignment \\
P2 scaling & Qwen2.5-\{0.5, 1.5, 3, 7, 14, 32\}B-Instruct & ✓ & C & varies & Scale-invariance curve \\
P2 ablation & \texttt{Qwen/Qwen2.5-Coder-7B-Instruct} (un-gated) & FC-trained & C & 4 & Tool-specialized variant cross-check \\
Legacy/Base & \texttt{NousResearch/Meta-Llama-3.1-8B}, \texttt{Mistral-7B-v0.3} (Base) & ✗ & — & — & Ablation only: "does K-bias work without FC training?" \\
\bottomrule
\end{tabularx}
\end{table}

\textbf{Important}: free-text scorers (Layer 1 of §5.2) apply to all models including Base; FC scorers (Layer 2–4) apply only to Instruct variants. Scaling curve and cross-family comparisons are Instruct-only for fair FC comparison. Our previously-run Llama-3.1-8B \textbf{Base} data (Wave 3a retry) is retained as "Base ablation" (§5.10 E10-b) only.


\textbf{Benchmarks.}


\begin{itemize}
\item \textbf{MetaTool Subtask1} (995 queries, 10 candidates + \texttt{None}; single-tool GT): scorer-invariance primary bed.
\item \textbf{MetaTool Subtask4} (497 queries, 2-tool GT): multi-tool + graded scoring primary bed. Ground-truth distribution: 100\% 2-tool.
\item \textbf{MMLU} (1000 samples, 5-shot): safety retention + hard-gate R-violation grid.
\item \textbf{WikiText-2} (full test, ctx=2048 non-overlap): compression (Thm 6.13 verification).
\item P3 stretch: BFCL-v3 Parallel, τ²-bench retail/airline, ToolAlpaca, HH-RLHF-500, ToxiGen-500.
\end{itemize}

\textbf{Steering hyperparameters.} Primary $\alpha=0.3$ (a0.2 is dead under strict scoring on all models). B\_ont built per (layer, KV-head) via Gram–Schmidt on catalog-derived facet sentences; rank $R=24$ for MetaTool ontology ($F=4$ facets). For Mistral: \texttt{skipL0 + pad-to-max} (validated fix, §5.13 E10).


\textbf{Evaluation invariants.} Greedy decoding; temperature 0; max\_new\_tokens 24 (single-tool) or 128 (multi-tool structured output); chat-template enabled for Instruct variants; function-calling via chat template's \texttt{tools} parameter with JSON tool schemas.


\subsection{5.2 Scoring framework (4-layer summary)}

A single forward pass per (method, model, query) emits predictions that are post-hoc scored under all applicable metrics. The layers:


\begin{table}[t]
\centering
\begin{tabularx}{\linewidth}{X | X | X}
\toprule
1. Free-text parsing & \texttt{substring\_any}, \texttt{first\_line}, \texttt{label\_logprob\{sum, mean\}} & Scorer-invariance triangulation on Subtask1 \\
\midrule
2. Function-calling & \texttt{fc\_name\_match}, \texttt{fc\_schema\_valid}, \texttt{fc\_label\_logprob} & Production-realistic (all scorers for Instruct models) \\
3. Set metrics & \texttt{F1}, \texttt{Jaccard}, \texttt{Exact-set}, \texttt{F\_\{0.5\}}, `EU($\alpha=1,\beta=2,\gamma=1$)` & Multi-tool symmetric + asymmetric cost (Subtask4) \\
4. Facet-graded & \texttt{FG-F1}, \texttt{FG-F\_\{0.5\}}, \texttt{FG-EU}, \texttt{ECE} (ambiguity subset) & Semantic proximity + calibration (Subtask4 + ambiguous Subtask1) \\
\bottomrule
\end{tabularx}
\end{table}

Layer-1–2 expose sensitivity of single-tool top-1 to parsing assumptions. Layer-3 captures multi-tool partial credit and production cost structure (wrong tool heavier than missing). Layer-4 credits same-facet-sibling predictions at $s=0.5$ and measures confidence calibration under ambiguity (Netsru Q8 alignment).


\textbf{Definitions.} Let $P$ = predicted tool multi-set, $G$ = ground-truth set, $\mathrm{TP} = |P \cap G|$, $\mathrm{FP} = |P \setminus G|$, $\mathrm{FN} = |G \setminus P|$.


\begin{itemize}
\item \textbf{F1}: $2 \cdot \mathrm{precision} \cdot \mathrm{recall} / (\mathrm{precision} + \mathrm{recall})$ — symmetric.
\item \textbf{F\_\{0.5\}}: $\frac{1.25 \cdot \mathrm{precision} \cdot \mathrm{recall}}{0.25 \cdot \mathrm{precision} + \mathrm{recall}}$ — precision weighted twice as heavily as recall (wrong tool heavier than missed).
\item \textbf{EU}: $\max(0,\; (\alpha\mathrm{TP} - \beta\mathrm{FP} - \gamma\mathrm{FN}) / (\alpha|G|))$ with $\alpha=1,\beta=2,\gamma=1$ — explicit cost model, clipped to [0,1].
\item \textbf{Jaccard}: $\mathrm{TP} / (\mathrm{TP} + \mathrm{FP} + \mathrm{FN})$.
\item \textbf{Exact-set}: $\mathbf{1}[P = G]$ (BFCL-v3 / τ²-bench default, strict).
\end{itemize}

Facet-graded similarity: each tool $t$ has facet tuple $\phi(t) = (\phi_1, \phi_2, \phi_3, \phi_4) \in$ (intent × domain × io\_type × category). Define $s(p, g) = 1$ if $p = g$, $0.5$ if $\exists f: \phi_f(p) = \phi_f(g)$, else $0$. \textbf{FG-F1} replaces $\mathrm{TP}$ with bipartite-matching $\sum s$; \textbf{FG-F\_\{0.5\}}, \textbf{FG-EU} analogously. \textbf{ECE} (expected calibration error) is computed per-query on the model's top-1 softmax probability vs. correctness on the ambiguity-flagged subset (ambiguous if more than one candidate shares $\phi$-dominant facet with GT).


\textbf{User-intuition validation} (ambiguous music query, GT=\{Spotify\}, 10 candidates):


\begin{table}[t]
\centering
\begin{tabularx}{\linewidth}{X | X | X | X | X}
\toprule
Prediction & F1 & F\_\{0.5\} & FG-F1 & Interpretation \\
\midrule
\{Spotify\} & 1.00 & 1.00 & 1.00 & exact \\
\{AppleMusic\} (same facet) & 0.00 & 0.00 & \textbf{0.50} & semantic neighbor credited under graded \\
\{Excel\} (cross-domain) & 0.00 & 0.00 & 0.00 & full penalty \\
\{Spotify, AppleMusic, YouTube\} & 0.40 & 0.45 & \textbf{0.75} & diffuse-in-facet, GT covered \\
\{AppleMusic, Excel\} & 0.00 & 0.00 & 0.17 & one neighbor + one cross-domain \\
\bottomrule
\end{tabularx}
\end{table}

This validates that $F_{0.5}$ and FG-F1 jointly encode the user requirements: wrong tool is heavy penalty (FP penalty 2× via F\_\{0.5\}), partial correctness rewarded (FG-F1 credits same-facet), cross-domain errors distinct from near-neighbor errors (s=0 vs s=0.5).


\subsection{5.3 Claim → experiment mapping (consolidated)}

Thirteen experiments partitioned across three priority tiers:


\textbf{Priority P1 (main paper, 91 GPU-hr)}:


\begin{itemize}
\item \textbf{E1} — Scorer-invariant mechanism specificity on Subtask1, 6 scorers × 3 B\_ont × 2 Instruct models (40 GPU-hr; Qwen partially done)
\item \textbf{E2} — Cor 6.9 decisive test on Subtask4, 9 metrics × 6 methods × 3 Instruct models (25 GPU-hr)
\item \textbf{E3} — Thm 6.1 per-sample bound: Qwen L13 + Llama L15, N=100 (15 GPU-hr, queued Wave 4)
\item \textbf{E4} — Cor 6.9 operator-level nrank SVD on 500 queries (2 GPU-hr)
\item \textbf{E5} — Rmk 6.14.A.3 R-violation MMLU grid, 25 cells (4 GPU-hr, queued R6)
\item \textbf{E6} — Thm 6.13 compression WT2 × \{Qwen, Llama\} × \{2, 3, 4\} bits (5 GPU-hr incremental)
\end{itemize}

\textbf{Priority P2 (reviewer defense + scaling, 60 GPU-hr)}:


\begin{itemize}
\item \textbf{E7} — Scaling curve Qwen2.5 \{0.5, 3, 7, 14\}B on Subtask4 FG-F1 (30 GPU-hr)
\item \textbf{E8} — Safety retention MMLU + HH-RLHF + ToxiGen (12 GPU-hr)
\item \textbf{E9} — Reproduced baselines CAA, ITI, PASTA, ASA, FocusDir, LoRA r=8, RAG (18 GPU-hr)
\item \textbf{E10} — Mistral closure (skipL0+padmax + Instruct H2) on Subtask4 (0 GPU-hr — Wave 3 ongoing)
\end{itemize}

\textbf{Priority P3 (future work, deferred)}:


\begin{itemize}
\item \textbf{E11} LoRA-R1 Thm 6.14 Hybrid · \textbf{E12} τ²-bench multi-turn · \textbf{E13} BFCL-v3 Parallel |G|-strat · \textbf{E14} zero-shot MetaTool→ToolAlpaca transfer · \textbf{E15} Thm 6.13 full bit curve · \textbf{E16} Conjecture 6.14 Full-FacetRot.
\end{itemize}

\textbf{Claim coverage (every theorem has a dedicated experiment):}


\begin{table}[t]
\centering
\begin{tabularx}{\linewidth}{X | X | X | X}
\toprule
Claim & Theorem/Remark & Primary Exp & Secondary \\
\midrule
C1 Geometric specificity (real≫random≫featshuffle) & — & E1 & E2 FG-F1 \\
C2 Phase-closure under Hypothesis (R) & Cor 6.7/6.8 & E5 & E3 \\
C3 ε-numerical-rank separation & Cor 6.9 & \textbf{E2} + E4 & E13 \\
C4 Categorical-channel compression & Thm 6.13 & E6 & E15 \\
C5 Attention-weighted bound & Thm 6.1 & E3 & — \\
C6 Hard-gate R violation & Rmk 6.14.A.3 & E5 & — \\
C7 Cross-model 2-family & — & E1 + E10 & E7 \\
C8 Scorer robustness & — & \textbf{E1} 6-scorer & E9 \\
C9 Ambiguity graded & §5.4.4 & \textbf{E2} FG-F1 gap & — \\
C10 Production alignment & Netsru Q8 & E2 FG-F\_\{0.5\}, EU + E8 & E12 \\
\bottomrule
\end{tabularx}
\end{table}

\textbf{Single-pass-multi-scorer design.} Each experiment emits all applicable scorer/metric outputs from one forward pass; no forward re-run is required for metric variants. This compresses the total cost from ~250 GPU-hr (naive) to ~150 GPU-hr (P1+P2).


\subsection{5.4 Results — E1 Scorer-invariant mechanism specificity (Subtask1, 995 queries)}

Subtask1 full 995 label\_logprob cross-model grid (Waves 1+2+3, complete 2026-04-15 02:30 KST):


\begin{table}[t]
\centering
\begin{tabularx}{\linewidth}{X | X | X | X | X | X | X | X}
\toprule
Model & Scorer & no\_steer & real a0.3 Δ & random a0.3 Δ & featshuffle a0.3 Δ & \textbf{real−random} & \textbf{real−featshuffle} \\
\midrule
Qwen2.5-7B-Instruct & substring\_any (legacy) & 75.58\% & +11.16pp & — & — & — & — \\
Qwen2.5-7B-Instruct & first\_line (parser-safe, codex Base) & 33.57\% & +2.81pp & −21.61pp & −32.16pp & \textbf{+24.42pp} & \textbf{+34.97pp} \\
Qwen2.5-7B-Instruct & label\_logprob \textbf{sum} & 52.46\% & +0.10pp & \textbf{−48.74pp} & \textbf{−40.10pp} & \textbf{+48.84pp} & \textbf{+40.20pp} \\
Qwen2.5-7B-Instruct & label\_logprob \textbf{mean} & 36.78\% & \textbf{+5.03pp} & \textbf{−23.01pp} & \textbf{−11.25pp} & \textbf{+28.04pp} & \textbf{+16.28pp} \\
Llama-3.1-8B-\textbf{Base} (NousResearch) & label\_logprob \textbf{sum} & 46.33\% & \textbf{+6.33pp} & −1.00pp & −0.20pp & \textbf{+7.33pp} & \textbf{+6.53pp} \\
Llama-3.1-8B-\textbf{Base} (NousResearch) & label\_logprob \textbf{mean} & 23.12\% & \textbf{+2.61pp} & −0.61pp & −1.41pp & \textbf{+3.22pp} & \textbf{+4.02pp} \\
Mistral-7B-v0.3 skipL0+padmax & label\_logprob \textbf{sum} & 69.35\% & \textbf{+3.12pp} & pending & pending & pending & pending \\
Mistral-7B-v0.3 skipL0+padmax & label\_logprob \textbf{mean} & 40.70\% & +0.20pp & pending & pending & pending & pending \\
Mistral-Instruct-v0.3 skipL0+padmax & label\_logprob \textbf{sum} & 61.51\% & \textbf{−2.92pp} & pending & pending & pending & pending \\
Mistral-Instruct-v0.3 skipL0+padmax & label\_logprob \textbf{mean} & 61.01\% & \textbf{−3.62pp} & pending & pending & pending & pending \\
Mistral-Instruct-v0.3 substring (Subtask1, full 995, 2026-04-15) & tool\_acc & 65.23\% & pending (real Mistral) & \textbf{+0.60pp} (random) & running (~19:30 KST) & pending & pending \\
\bottomrule
\end{tabularx}
\end{table}

\textbf{Cross-model 3-family positive under strict label\_logprob (Qwen + Llama + Mistral-Base)}: Qwen sum +0.10 / mean +5.03, Llama-Base sum +6.33 / mean +2.61, Mistral-Base-v0.3 (skipL0+padmax fix) sum +3.12 / mean +0.20. All three base architecture families register positive. Mistral-\textbf{Instruct}-v0.3 is the sole negative (sum −2.92, mean −3.62): the Instruct variant's no\_steer is itself 7.84pp \textbf{below} the Base variant (61.51\% vs 69.35\%), and K-bias further degrades — consistent with chat-template hedging rather than a mechanism counterexample (analysis §5.5.1).


\textbf{Null-control specificity (Qwen mean sharpest, complete 4-cell)}: real +5.03 vs random −23.01 vs featshuffle −11.25 → gaps +28.04 / +16.28. \textbf{Direction specificity is scorer-invariant and model-invariant}: ordering real ≫ featshuffle ≥ random holds wherever the full 3-control triple is populated (Qwen sum/mean, Llama sum/mean, codex first\_line).


\textbf{Headline accuracy is scorer-dependent} (+0.1 to +11.15pp). \textbf{Mechanism specificity is scorer-invariant}: under every strict scorer, the ordering real > random > featshuffle holds with gaps +16 to +49pp — between one and two orders of magnitude larger than the accuracy headline. The "any projector works" alternative hypothesis is decisively rejected.


\textbf{Answerability vs discrimination decomposition} (under codex first\_line parser\_safe, full 995):


\begin{itemize}
\item Original a0.3: matched-rate +2.81pp, conditional-accuracy +1.37pp → small real discrimination.
\item Opaque a0.3: matched-rate +20.30pp, conditional-accuracy −4.44pp → \textbf{new-commit correctness 65.82\% (6.6× random)} → the "answerability rescue" IS semantic routing, not artifact (§5.4.1 analysis).
\end{itemize}

Llama-3.1-8B Base full 3-control is complete (row 5–6 above): sum real +6.33 / random −1.00 / featshuffle −0.20 (gap +7.33 / +6.53); mean real +2.61 / random −0.61 / featshuffle −1.41 (gap +3.22 / +4.02). Second family triple verified.


\subsubsection{5.4.1 Subtask1 Q-coverage and K-bias single-tool accuracy lift — \textbf{cross-model under matched scorer}}

\textbf{Critical reviewer clarification (2026-04-15)}: The Qwen sum label-logprob +0.10pp cell in the main §5.4 table is under the \emph{strictest} closed-set scorer. A direct same-scorer comparison with Llama requires reading the substring-scorer row for both models, which we now align:


\begin{table}[t]
\centering
\begin{tabularx}{\linewidth}{X | X | X | X | X | X}
\toprule
Model & Scorer & no\_steer & K-bias α=0.3 & Q-cov β=−0.3 & Q-cov β=−0.1 \\
\midrule
Qwen-Inst Subtask1 & substring\_any (legacy) & 75.58\% & \textbf{+11.16pp} (legacy memory, full 995) & pending & pending \\
Qwen-Inst Subtask1 & label\_logprob sum (strict) & 52.46\% & +0.10pp & — & — \\
Qwen-Inst Subtask1 & label\_logprob mean (strict) & 36.78\% & +5.03pp & — & — \\
Qwen-Inst Subtask1 & substring / tool\_acc (this paper) & 60.30\% & \textbf{+1.41pp} (2026-04-15) & +4.12pp & +3.22pp \\
\textbf{Llama-Inst Subtask1} & \textbf{substring / tool\_acc} & \textbf{62.31\%} & \textbf{+15.08pp ⚡} & \textbf{+8.04pp} & +0.30pp \\
\bottomrule
\end{tabularx}
\end{table}

\textbf{Same-scorer cross-model comparison} (substring scorer, strict tool-candidate match):


\begin{itemize}
\item Qwen2.5-7B-Instruct: +1.41pp
\item Llama-3.1-8B-Instruct: \textbf{+15.08pp}
\end{itemize}

\subsubsection{5.4.1.1 Cross-model K-bias lift asymmetry — empirical observation and refuted mechanism (revised 2026-04-15)}

The substring-scorer K-bias lift ratio $+15.08 / +1.41 \approx 10.7\times$ (Llama-Inst > Qwen-Inst on Subtask1 full 995) is a reliable empirical observation. Our initial theoretical decomposition via Thm 6.1 bound factors was subsequently \emph{falsified by direct measurement}; this section reports the refutation honestly and scopes the claim to the empirical observation with correct-mechanism pointers.


\textbf{Initial (falsified) decomposition.} We hypothesized three architectural factors jointly predict the 10.7× ratio:


\begin{itemize}
\item (a) Mode A (Llama) vs Mode C (Qwen): attention-weighted V-variance $\mathrm{Var}_s V$ 5–10× larger on Mode A.
\item (b) GQA 4:1 vs 7:1: effective per-Q-head K-bias magnitude 1.75× larger on Llama.
\item (c) $B_{\mathrm{ont}}$ rank 19 vs 24: 1.26× in the opposite direction.
\end{itemize}

Predicted ratio $(5\text{–}10) \times 1.75 / 1.26 \approx 7\text{–}14\times$, matching 10.7× observation "within range".


\textbf{Refutation (direct Thm 6.1 verification, N=100 queries × 32 heads, 2026-04-15)}. We measured the Thm 6.1 decomposition on Llama-3.1-8B-Instruct at layer 15 (\texttt{reports/thm61\_llama\_2026\_04\_15/llama\_L15\_a0.3\_N100.json}) and compared to Qwen2.5-7B-Instruct at layer 13 (\texttt{reports/theory\_verify\_2026\_04\_14/thm61\_qwen\_L13\_a0.3\_N100.json}):


\begin{table}[t]
\centering
\begin{tabularx}{\linewidth}{X | X | X | X}
\toprule
Factor (per-head mean at $\alpha_K = 0.3$) & Qwen L=13 & Llama L=15 & Qwen/Llama \\
\midrule
$\mathbb E[\|\hat o - o\|^2]$ (LHS) & 0.509 & 0.059 & \textbf{8.69×} \\
$\mathbb E[\mathrm{qaMSE} \cdot \mathrm{Var}_s V]$ (RHS leading) & 19.73 & 2.41 & \textbf{8.18×} \\
qaMSE & 0.207 & 0.239 & 0.87× \\
$\mathrm{Var}_s V$ (raw) & 49.16 & 6.85 & \textbf{7.17×} \\
$V_{\max}$ (per-head max V norm) & 12.94 & 5.66 & \textbf{2.29×} \\
$\mathrm{Var}_s V / V_{\max}^2$ (scale-normalized) & 0.308 & 0.218 & \textbf{1.42×} \\
\bottomrule
\end{tabularx}
\end{table}

The factor (a) 5–10× prediction is empirically falsified \emph{in the wrong direction}: Qwen $\mathrm{Var}_s V$ is 7.17× \emph{larger} than Llama, dominated by V-scale (2.29² = 5.24× of the 7.17× ratio) with only 1.42× remaining after scale-normalization. The Thm 6.1 bound magnitude (LHS, RHS-leading) is also Qwen-larger by 8.18×, opposite the 10.7× Llama-larger K-bias lift direction.


\textbf{Mechanistic clarification (why bound ratio ≠ lift ratio)}. The Thm 6.1 bound $\|\hat o - o\|^2 \le 2\,\mathrm{qaMSE} \cdot \mathrm{Var}_s V + C_1\rho^4$ characterizes the \emph{magnitude of attention-output perturbation error}, not the accuracy lift. The K-bias accuracy lift is governed by the first-order coefficient $G_K = \langle \nabla_{\Delta_K} \log p, B_{\mathrm{ont}} B_{\mathrm{ont}}^\top K\rangle$ of Thm 6.17 (iii), which measures gradient alignment of the K-bias direction with the log-probability increase, a distinct dimensionless quantity. Cross-model bound magnitude ratios scale with model weight scales (V norm, K norm); cross-model lift ratios scale with gradient alignment and saturation dynamics. These can diverge, as our measurement shows (+8× bound ratio in one direction, +10.7× lift ratio in the opposite direction).


**Direct $G_K$ measurement (partial, 2026-04-15)**. We measured $|\cos(\nabla_{\Delta_K} \log p, B_{\mathrm{ont}} B_{\mathrm{ont}}^\top K)|$ per head on 50 Subtask4 queries (\texttt{reports/gk\_measurement\_2026\_04\_15/}):


\begin{table}[t]
\centering
\begin{tabularx}{\linewidth}{X | X | X | X}
\toprule
Quantity & Qwen L=13 & Llama L=15 & Llama/Qwen \\
\midrule
$\|\cos(\nabla_{\Delta_K} \log p, \text{direction})\|$ & 0.098 & 0.162 & \textbf{1.66×} \\
\bottomrule
\end{tabularx}
\end{table}

The cosine alignment ratio 1.66× is the dimensionless, weight-scale-invariant gradient measurement. It captures approximately 1/6 of the observed 10.7× lift ratio. The remaining factor of $\sim 6\text{–}7\times$ plausibly reflects (i) saturation dynamics (Mode A's lift grows near-linearly up to $\alpha = 0.3$ while Mode C saturates early), (ii) the 1.75× GQA group-size factor, and (iii) task-specific first-token gradient effects from tokenizer-level tool-name fragmentation differences. Full characterization of the $\alpha$-saturation curve and gradient decomposition per model is queued as future work.


\textbf{Scope update (honest)}. We report the 10.7× ratio as an empirical observation with the following confirmed properties: (i) direction-specificity ordering (real ≫ featshuffle ≥ random) holds cross-model under every tested scorer; (ii) gradient alignment is 1.66× Llama-larger (partial mechanism); (iii) the Thm 6.1 bound magnitude goes in the \emph{opposite} direction, reflecting bound-magnitude characterization of \emph{error} distinct from lift's \emph{alignment}. The original 7–14× decomposition is retracted. A scorer-robustness check of the 10.7× ratio on Llama-Inst under label-logprob (currently in progress) will further scope the observation.


\textbf{Llama-Base vs Llama-Inst asymmetry (preserved)}. GQA-4, Mode A, identical head\_dim and B\_ont rank do not change between Base and Instruct variants, yet K-bias lift is 2.4× larger on Instruct (+6.33pp Base → +15.08pp Inst). This within-architecture asymmetry is an independent observation not explained by our refuted Thm 6.1 decomposition; we interpret it as SFT-recipe-induced shift in either gradient alignment $G_K$ or saturation behavior. Direct measurement at matched SFT recipe (e.g., re-SFT Llama-Base on Llama-Inst's instruction corpus) is queued as future work.


\textbf{Falsifiability (remaining)}: if $G_K$ per-model measurement under $\alpha \to 0$ fully accounts for the lift ratio (including the $\sim 6\times$ residual), the mechanism is gradient-alignment-dominated. If not, saturation dynamics (lift-vs-$\alpha$ curve cross-model) becomes the dominant factor; scaling-curve prediction on Qwen2.5-3B / Llama-3.1-70B is the discriminator (§5.11).


Beyond the label\_logprob cells of the above table,


Beyond the label-logprob cells of the above table, we also evaluated Q-coverage and K-bias under the legacy substring scorer at full 995 (PM Wave 2 results, 2026-04-15):


\begin{table}[t]
\centering
\begin{tabularx}{\linewidth}{X | X | X | X}
\toprule
Method & top1 & Δ vs no\_steer 60.30\% & preds (matched / no\_match / none) \\
\midrule
no\_steer & 60.30\% & — & 821 / 43 / 131 \\
\textbf{ocq\_qbias\_b−0.3} & \textbf{64.42\%} & \textbf{+4.12pp} ★ & 853 / 28 / 114 \\
\textbf{ocq\_qbias\_b−0.1} & \textbf{63.52\%} & \textbf{+3.22pp} & 860 / 17 / 118 \\
\textbf{ocq\_bias\_a0.3 (K-bias)} & \textbf{61.71\%} & \textbf{+1.41pp} ✅ & 825 / 71 / 99 \\
\bottomrule
\end{tabularx}
\end{table}

\textbf{Key findings}:


\begin{enumerate}
\item \textbf{K-bias produces single-tool accuracy lift} (+1.41pp full 995). This contradicts the §5.5.2 multi-tool failure (−4.6pp) — K-bias works as Cor 6.9 originally predicted on \emph{single-tool} tasks where no autoregressive coverage challenge arises, and only fails on multi-tool emission. K-channel is therefore \emph{not} "stability-only" as our prior re-scope suggested; it has a verified single-tool accuracy contribution.
\item \textbf{Q-coverage cross-task universality}: positive on both Subtask1 (+3.22 to +4.12pp) and Subtask4 (+1.6pp).
\item \textbf{β-optimum is task-dependent}: Subtask4 prefers gentler β=−0.1 (multi-tool needs careful coverage), Subtask1 tolerates aggressive β=−0.3 (single-tool benefits from sharper attention reallocation).
\end{enumerate}

\subsection{5.5 Results — E2 Cor 6.9.6 stability characterization (Subtask4, 497 × 2-tool)}

\textbf{Stability rather than accuracy.} Cor 6.9 was originally used to predict a \emph{multi-tool accuracy lift} on the hypothesis that rank-$R$ support would enable simultaneous emission of $R$-facet-aligned tool names in a single attention pass (call this the "F-simultaneous accuracy" hypothesis). Full-scale measurement falsifies this prediction: real $B_{\mathrm{ont}}$ $\alpha=0.3$ F1 = 0.685 vs no\_steer 0.731, $\Delta = -4.6$pp. Autoregressive re-attention (§5.5 discussion below) prevents a \emph{stationary} K-bias from driving facet-wise coverage across decoding steps, regardless of operator spectral rank. The originally-predicted multi-tool accuracy lift requires a non-stationary K-bias (§5.5.2, Thm 6.9.5/6.15).


However, the \emph{same} rank-$R$ operator structure manifests as a \textbf{stability property} whose empirical signal is an order of magnitude stronger than the originally-predicted accuracy lift. The following subsection verifies Corollary 6.9.6 (§3.3) at full scale on Subtask4.


\textbf{Planned experimental cells} (launch queued post Wave 4):


\begin{table}[t]
\centering
\begin{tabularx}{\linewidth}{X | X | X}
\toprule
no\_steer & Qwen-Instruct, Llama-Instruct, Mistral-Instruct & F1, F\_\{0.5\}, EU, Jaccard, Exact, FG-F1, FG-F\_\{0.5\}, FG-EU, ECE \\
\midrule
a0.3 real & same 3 & same 9 metrics \\
a0.3 random & same 3 & same \\
a0.3 featshuffle & same 3 & same \\
AdaSEKA 2-expert & same 3 & same \\
AdaSEKA 3-expert & same 3 & same \\
\bottomrule
\end{tabularx}
\end{table}

Total: 18 forward-pass configurations × 9 metrics = 162 numbers. Expected runtime 25 GPU-hr.


\textbf{Theorem-level prediction (Cor 6.9)}: for any max-normalized-routing baseline, recall on 2-tool queries is capped at 0.5 by construction (one expert → one tool emission). Therefore $\mathrm{F_{0.5}} \le \tfrac{1.25 \cdot 1 \cdot 0.5}{0.25 \cdot 1 + 0.5} \approx 0.83$. Our facet-gated method has no such cap (rank $R=24$ supports F-simultaneous emission); $\mathrm{F_{0.5}}$ up to 1.0 achievable. \textbf{This is a falsifiable numerical prediction}.


\textbf{Subtask4 N=20 smoke complete (2026-04-15 00:45 KST, Qwen2.5-7B-Instruct, all 3 B\_ont variants)}:


\begin{table}[t]
\centering
\begin{tabularx}{\linewidth}{X | X | X | X | X | X | X | X}
\toprule
B\_ont & Method & F1 & F\_0.5 & EU & Jaccard & Exact & Recall \\
\midrule
real & no\_steer & 0.550 & 0.550 & 0.300 & 0.467 & 0.300 & 0.550 \\
real & a0.3 & \textbf{0.533} & 0.542 & 0.150 & 0.408 & 0.150 & 0.525 \\
random & no\_steer & 0.550 & 0.550 & 0.300 & 0.467 & 0.300 & 0.550 \\
random & a0.3 & \textbf{0.000} & 0.000 & 0.000 & 0.000 & 0.000 & 0.000 \\
featshuffle & no\_steer & 0.550 & 0.550 & 0.300 & 0.467 & 0.300 & 0.550 \\
featshuffle & a0.3 & \textbf{0.000} & 0.000 & 0.000 & 0.000 & 0.000 & 0.000 \\
\bottomrule
\end{tabularx}
\end{table}

\textbf{Gap (real − random / real − featshuffle) = +53.3pp on F1} (and all other metrics) — decisive mechanism-specificity at N=20.


\textbf{Reinterpretation of Cor 6.9 on Subtask4}: The predicted signature "rank R supports multi-tool emission accuracy lift" is not observed: real a0.3 F1 ≈ no\_steer F1 (no accuracy improvement). However, the \textbf{null-control collapse} observed (random/featshuffle both produce F1 = 0.000 — the model emits no parseable \texttt{<tool\_call>} blocks under random/featshuffle K-bias at α=0.3) reveals a stronger empirical signature: \textbf{the ontology subspace is the unique α=0.3-magnitude K-perturbation direction that preserves the model's structured-output emission capability}. Random and feature-shuffle perturbations of matched magnitude destroy the chat-template FC generation completely.


We reformulate the Cor 6.9 downstream signature:


\begin{quote}
\textbf{Geometric-safety interpretation.} For any FC-trained instruction model, there is a characteristic K-perturbation magnitude α* above which arbitrary-direction K-biases break structured output. The facet-gated operator's rank-$R$ ontology subspace is the unique direction that remains within the model's "natural tool-reasoning manifold" at α up to at least 0.3; other directions of the same magnitude exit that manifold and collapse emission. This is a \emph{stability} (not an \emph{accuracy-improvement}) manifestation of Cor 6.9's rank structure.
\end{quote}

This reinterpretation is consistent with:


\begin{itemize}
\item Section 5.7 E4 (operator-level nrank: ours = 24, AdaSEKA = 6–8 depending on T): the rank gap exists as predicted.
\item Section 5.8 E5 hard-gate R-violation (MMLU): discontinuous gates that violate Lipschitz regularity degrade monotonically in α — same "safe-direction" intuition on a different benchmark.
\end{itemize}

\textbf{Autoregressive re-attention limitation (§5.5.1)}: the original "F-simultaneous rank R → multi-tool emission" prediction assumed that steering toward a multi-facet K subspace would cause the model to emit multiple tool\_calls per query. In practice, multi-tool emission relies on \emph{sequential} attention re-computation across decoding steps (context updates alter query direction, enabling coverage of un-emitted facets). Time-invariant K-bias applied uniformly across steps does not compose with this sequential mechanism: boosting facet-aligned attention equally at every step does not drive the model toward complementary facets in later steps.


Proposed fix under investigation (§5.11 future work E11'): a KQV-hybrid where (i) K-bias marks facet structure (small α\_K), (ii) V-bias amplifies in-ontology V content (α\_V moderate), and (iii) Q-side coverage-masked projection removes emitted-facet direction from the query at each step. Theorem 6.15 (proposed, Appendix B.7.8.1) formalizes this combination. V-bias smoke under way (§5.12 live run).


\textbf{Full 497 Subtask4 results (2026-04-15 02:30 KST, all 3 B\_ont variants complete)}:


\begin{table}[t]
\centering
\begin{tabularx}{\linewidth}{X | X | X | X | X}
\toprule
B\_ont & Method & F1 & Recall & Exact \\
\midrule
real & no\_steer & \textbf{0.731} & 0.716 & 0.525 \\
real & a0.3 & 0.685 & 0.672 & 0.473 \\
random & no\_steer & 0.731 & 0.716 & 0.525 \\
random & a0.3 & \textbf{0.000} & 0.000 & 0.000 \\
featshuffle & no\_steer & 0.731 & 0.716 & 0.525 \\
featshuffle & a0.3 & \textbf{0.000} & 0.000 & 0.000 \\
\bottomrule
\end{tabularx}
\end{table}

\textbf{Real − random gap = real − featshuffle gap = +68.5pp F1} at full N=497. The null-control collapse hypothesized from smoke (§5.5) is \textbf{decisively verified at full scale}: α=0.3 random and featshuffle K-bias completely destroy structured \texttt{<tool\_call>} emission, while the ontology direction preserves F1 within 4.6pp of no\_steer. This is the strongest direction-specificity evidence in the paper.


\textbf{Final interpretation of Cor 6.9 downstream signature (empirical verdict)}:


\begin{itemize}
\item \textbf{Accuracy-lift version (original prediction): FALSIFIED}. Real a0.3 F1 ≤ no\_steer F1 on both smoke (N=20, Δ=−1.7pp) and full (N=497, Δ=−4.6pp).
\item \textbf{Geometric-safety version (reframed): VERIFIED on smoke; full 497 null-control pending for confirmation}. Random/featshuffle at α=0.3 produced F1=0.000 on N=20 — complete collapse vs real's preserved 0.53. Expected to hold on full 497.
\end{itemize}

\textbf{Paper claim for Subtask4 (final, full-scale verified)}:


\begin{quote}
"Cor 6.9 predicts the ontology direction is the unique α=0.3-magnitude K-perturbation that preserves FC-structured-output emission on multi-tool queries. Empirically (Qwen2.5-7B-Instruct, MetaTool Subtask4 \textbf{full 497}): real a0.3 maintains F1=0.685 (no\_steer 0.731, Δ=−4.6pp), while \textbf{random/featshuffle both collapse to F1=0.000} — a +68.5pp direction-specificity gap. This is a \emph{stability} manifestation of the rank separation consistent with Cor 6.9's operator-level rank bound (§5.7 E4: 24.0 vs 7.44), distinct from the originally predicted accuracy lift. Multi-tool emission under stationary K-bias is limited by autoregressive re-attention; Thm 6.15 (KQV hybrid, App. B.7.8.1) proposes a theoretically-motivated fix, and §5.5.2 reports a first empirical improvement via contrastive K-bias (Thm 6.9.5 family)."
\end{quote}

\subsection{5.5.2 Non-uniform K-bias extension — smoke positive, full 497 NEGATIVE (artifact)}

The §5.5 stability result is the \emph{baseline} operating point of the facet-gated operator. Accuracy lift on multi-tool queries requires a \textbf{non-stationary} K-bias that evolves across decoding steps (Thm 6.9.5/6.15, Appendix B.7.8). The contrastive subfamily (per-step subtraction of dominant-singular-direction K content) was the simplest non-stationary instantiation. We report both the original positive smoke signal and the \emph{negative full-scale replication} below; the headline is that the smoke signal does \textbf{not} survive at full scale.


\textbf{Independence from §5.5.} This subsection's negative result does \emph{not} affect §5.5's stability claim. Cor 6.9.6's $+68.5$pp direction-specificity gap on full 497 is independent of any extension's empirical fate.


14-configuration sweep at reduced magnitude and with contrastive / normalized variants (smoke N=20 MetaTool Subtask4, 2026-04-15 02:30 KST):


\begin{table}[t]
\centering
\begin{tabularx}{\linewidth}{X | X | X | X}
\toprule
Variant & α / params & F1 (vs no\_steer 0.550) & Δ \\
\midrule
flat real (baseline reference) & a=0.3 & 0.533 & −0.017 \\
α-sweep & a=0.05 & 0.550 & 0.000 \\
α-sweep & a=0.10 & 0.533 & −0.017 \\
α-sweep & \textbf{a=0.15} & \textbf{0.575} & \textbf{+0.025} \\
α-sweep & a=0.20 & 0.492 & −0.058 \\
normalized (Thm 6.9.5) & a=0.1 & 0.500 & −0.050 \\
normalized & a=0.3 & 0.325 & −0.225 \\
normalized & a=0.5 & 0.000 & −0.550 \\
normalized & a=1.0 & 0.025 & −0.525 \\
contrastive & a=0.3 d=1 & 0.583 & +0.033 \\
contrastive & a=0.3 d=2 & 0.508 & −0.042 \\
\textbf{contrastive} & \textbf{a=0.3 d=3} & \textbf{0.608} & \textbf{+0.058} \\
contrastive & a=0.5 d=1 & 0.067 & −0.483 \\
contrastive & a=0.5 d=2 & 0.225 & −0.325 \\
contrastive & a=0.5 d=3 & 0.067 & −0.483 \\
\bottomrule
\end{tabularx}
\end{table}

\textbf{Key finding}: at α=0.3, contrastive depth-3 K-bias yields \textbf{F1 = 0.608 (+5.8pp over no\_steer 0.550)} on the same smoke set where flat real a0.3 gives 0.533 (−1.7pp). This is the \textbf{first positive Subtask4 F1 signature} for any member of the K-bias family, and it is predicted by Thm 6.9.5/6.15 (non-uniform family): contrastive mixing injects facet direction while subtracting paired sibling-facet leakage, directly targeting the autoregressive re-attention limitation identified in §5.5. \textbf{V-bias alone fails} (max F1 0.558 at ak=0.1/av=0.1, no config beats flat real). \textbf{Normalized-only variant fails at α ≥ 0.3}: catastrophic collapse, matching the under-gating regime of Cor 6.9.4.


\textbf{Full 497 contrastive verification (2026-04-15 09:25 KST) — smoke signal does NOT replicate}:


\begin{table}[t]
\centering
\begin{tabularx}{\linewidth}{X | X | X | X}
\toprule
Method & smoke F1 (N=20) & full F1 (N=497) & Δ (full − no\_steer 0.731) \\
\midrule
no\_steer & 0.550 & 0.731 & — \\
ocq\_cbias\_a0.3\_d1 & 0.583 (+3.3pp) & 0.690 & \textbf{−4.1pp} \\
ocq\_cbias\_a0.3\_d3 & 0.608 (+5.8pp) & 0.695 & \textbf{−3.6pp} \\
\bottomrule
\end{tabularx}
\end{table}

The +5.8pp smoke signal does not survive full-scale replication (full Δ = −3.6pp on d=3, similar on d=1). We classify the smoke result as a \emph{small-N variance artifact} (N=20 has bootstrap standard error ≈ 0.11 on F1; the +5.8pp signal is within one-sigma of per-sample variance on this benchmark). Stationary K-bias of \emph{any} form tested — flat, normalized (Thm 6.9.5 literal), or contrastive — cannot drive multi-tool coverage at full scale; this corroborates §5.5's reframing that the autoregressive re-attention barrier is structural for stationary perturbations.


\textbf{Live promising signal (Q-coverage subtraction, smoke N=20)}:


\begin{table}[t]
\centering
\begin{tabularx}{\linewidth}{X | X | X}
\toprule
Method & F1 (smoke N=20) & Δ vs no\_steer 0.550 \\
\midrule
ocq\_qbias\_b−0.1 (Q-coverage subtract, weak) & \textbf{0.658} & \textbf{+10.8pp} \\
ocq\_qbias\_b−0.3 & 0.575 & +2.5pp \\
ocq\_qbias\_b−0.5 & 0.600 & +5.0pp \\
ocq\_qkv\_a0.3\_v0\_q−0.3 (K + Q-coverage) & 0.500 & −5.0pp \\
ocq\_qkv\_a0.3\_v0.3\_q−0.3 (full QKV joint Thm 6.17) & 0.500 & −5.0pp \\
\bottomrule
\end{tabularx}
\end{table}

Q-only coverage subtraction at $\beta = -0.1$ is the strongest non-stability multi-tool signal observed in this paper (smoke +10.8pp). It is the \textbf{isolated Thm 6.17 (b) component} (Q-coverage gradient direction), without K-marker or V-amplifier. Crucially, \emph{adding} the K-marker at $\alpha_K=0.3$ destroys the Q-only gain — they interact destructively at this magnitude rather than additively as the first-order Lagrangian decomposition (Thm 6.17 (d)) predicted.


\textbf{Full-497 verification of Q-coverage β=−0.1 (Qwen2.5-7B-Instruct, complete 2026-04-15 12:18 KST)}:


\begin{table}[t]
\centering
\begin{tabularx}{\linewidth}{X | X | X | X | X | X}
\toprule
Method & F1 & F0.5 & Recall & Exact & Δ vs no\_steer 0.731 \\
\midrule
no\_steer & 0.731 & 0.745 & 0.716 & 0.525 & — \\
\textbf{ocq\_qbias\_b−0.1 (real B\_ont)} & \textbf{0.747} & \textbf{0.763} & \textbf{0.763} & 0.527 & \textbf{+1.6pp F1, +4.7pp recall} ✅ \\
\bottomrule
\end{tabularx}
\end{table}

The smoke +10.8pp shrinks to +1.6pp at full scale (small-N variance regression), but the \emph{positive sign and recall lift} survive. F0.5 lifts +1.8pp, recall +4.7pp, Exact +0.2pp — the recall channel is where the signal lives, consistent with Q-coverage driving multi-tool emission of a previously-unsaid second tool.


\textbf{β-sweep at full 497 confirms refined Thm 6.17′ small-α regime}:


\begin{table}[t]
\centering
\begin{tabularx}{\linewidth}{X | X | X}
\toprule
β & F1 & Exact \\
\midrule
0 (no\_steer) & 0.731 & 0.525 \\
−0.05 & 0.730 & 0.533 \\
\textbf{−0.1} & \textbf{0.747} ★ & 0.527 \\
−0.15 & 0.729 & 0.499 \\
−0.2 & 0.727 & 0.493 \\
−0.3 & 0.622 & — \\
−0.5 & 0.614 & — \\
−0.7 & 0.000 & — \\
\bottomrule
\end{tabularx}
\end{table}

\textbf{Single isolated peak at β=−0.1}, ±0.05 outside loses lift. The empirical $\alpha_{\mathrm{coupling}} \approx 0.1$ value of refined Thm 6.17′ is measured to ±0.05 precision.


\textbf{Null-control falsifiability — Q-coverage is ontology-specific (decisive Thm 6.17 verification, complete 2026-04-15 12:42 KST)}:


\begin{table}[t]
\centering
\begin{tabularx}{\linewidth}{X | X | X | X | X}
\toprule
B\_ont source & Method & F1 & Δ vs no\_steer 0.731 & Interpretation \\
\midrule
\textbf{real (ontology)} & ocq\_qbias\_b−0.1 & \textbf{0.747} & \textbf{+1.6pp} ✅ & unique lift \\
\textbf{featshuffle} & ocq\_qbias\_b−0.1 & 0.725 & −0.6pp & structure-preserving null fails \\
\textbf{random} & ocq\_qbias\_b−0.1 & 0.707 & −2.4pp & noise null fails \\
\bottomrule
\end{tabularx}
\end{table}

Real − featshuffle gap = \textbf{+2.2pp F1}. Real − random gap = \textbf{+4.0pp F1}. The three-tier ordering (real ≫ structure-preserving null ≫ noise null) precisely matches the Thm 6.17 (b) prediction: the gradient $\nabla_{\Delta_Q} \log p$ has support on the \emph{ontology} subspace, \emph{not} on arbitrary rank-24 directions and \emph{not} even on subspaces that preserve channel-marginal statistics. Featshuffle preserves per-channel norms and variances (only permutes feature indices in $B_{\mathrm{ont}}$) yet still fails to lift — this is the strongest formulation of the falsifiability claim possible.


This null-control result is the \emph{second independent ontology-specificity verification} of the paper, after the §5.5 stability gap (+68.5pp F1) on the same B\_ont. The two-channel verification (stability + accuracy lift) jointly forecloses the strongest reviewer counter-hypothesis ("rank-24 with arbitrary basis would suffice") at decisive p-value.


\textbf{Cross-model verification on Llama-3.1-8B-Instruct (Subtask4 full 497)}:


\begin{table}[t]
\centering
\begin{tabularx}{\linewidth}{X | X | X}
\toprule
Method & Llama-Inst F1 & Δ vs no\_steer 0.623 \\
\midrule
no\_steer & 0.623 & — \\
K-bias α=0.3 & 0.311 & \textbf{−31.2pp} (Llama α* < 0.3, FC manifold violated) \\
\textbf{Q-coverage β=−0.1} & \textbf{0.627} & \textbf{+0.4pp} ✅ \\
\bottomrule
\end{tabularx}
\end{table}

Q-coverage's lift is smaller on Llama-Inst than on Qwen-Inst (+0.4 vs +1.6pp) but the sign and the safety property survive. Crucially, K-bias at the same magnitude $\alpha=0.3$ catastrophically collapses Llama (−31.2pp) while Q-coverage at $\beta=−0.1$ remains in-manifold. \textbf{Q-coverage is therefore the universally-safe member of the perturbation family} — Qwen tolerates K-bias at α=0.3, Llama does not, and only Q-coverage at β=−0.1 stays on-manifold across both.


\textbf{Cross-model channel asymmetry (2026-04-15, honest acknowledgment)}. The 4× Qwen/Llama ratio in Q-coverage lift (Qwen +1.64 vs Llama +0.4) contrasts directly with the cross-model K-bias lift ordering, where \emph{Llama dominates Qwen} (Subtask1 K-bias +15.08 Llama vs +1.41 Qwen at substring scorer, ratio 10.7×). Neither channel uniformly favors one model:


\begin{table}[t]
\centering
\begin{tabularx}{\linewidth}{X | X | X | X}
\toprule
Channel & Qwen Subtask4 Δ & Llama Subtask4 Δ & Qwen/Llama ratio \\
\midrule
Q-coverage ($\beta=-0.1$) & +1.64pp & +0.40pp & \textbf{4.1× (Qwen stronger)} \\
K-bias ($\alpha=0.3$, Subtask1) & +1.41pp & +15.08pp & \textbf{0.09× (Llama stronger)} \\
\bottomrule
\end{tabularx}
\end{table}

We interpret this as \textbf{complementary channel design}: Q-coverage is more effective on Mode-C (concentrated) models where step-adaptive subtraction of already-emitted facets has clear attention mass to redistribute, while K-bias is more effective on Mode-A (diffuse) models where there is room to concentrate attention toward facet-aligned keys. The combined dual-channel operator (Q-coverage + small-α K augmentation of §3.6.1 (iii)) is model-adaptive: Qwen's strength is primarily the Q-channel, Llama's is primarily the K-channel, and the verified pair (Q+K at $\alpha_K = 0.025$, $\beta_Q = -0.1$ on Subtask4) captures both operating regimes. This reframes the cross-model asymmetry as \emph{evidence for designing the two-channel operator} rather than as a generalization weakness of either channel alone.


\textbf{Cross-task confound caveat (honest)}. The complementary-channel inversion as observed above mixes \emph{task} with \emph{model}: the Qwen-stronger result is on Subtask4 (Q-only), the Llama-stronger result is on Subtask1 (K-only). A clean 2×2 (model × task) matrix would require also Llama Subtask4 K-only and Qwen Subtask1 Q-only at matched scorer, with at least one channel showing model-preference inversion \emph{within the same task} to support a model-Mode-driven (not task-driven) interpretation. We have one half of this matrix:


\begin{table}[t]
\centering
\begin{tabularx}{\linewidth}{X | X | X | X | X}
\toprule
 & Qwen Subtask4 & Llama Subtask4 & Qwen Subtask1 & Llama Subtask1 \\
\midrule
Q-only β=-0.1 & +1.64pp ✅ & +0.40pp ✅ & +3.22pp (substring) & (queued) \\
K-only α=0.3 & −4.6pp (stability) & −31.2pp (catastrophic) & +1.41pp & +15.08pp ✅ \\
\bottomrule
\end{tabularx}
\end{table}

The decisive missing cell is \textbf{Llama Subtask4 K-only α sweep} (small-α regime, $\alpha \in \{0.05, 0.1, 0.2, 0.3\}$). If Llama Subtask4 small-α K-only shows positive lift dominating Q-only's +0.40pp on the same task, the model-Mode interpretation is supported within Subtask4. If small-α K-only is similarly small or negative on Llama Subtask4, the task-driven (not model-driven) interpretation is more parsimonious. This 2×2-completion experiment is queued (~3 GPU-hr on A6000); we present the complementary-channel reading as a \emph{working hypothesis} rather than as fully verified, with the falsifiability path explicitly named. Until the Llama Subtask4 K sweep lands, the paper's central operational claim stays at the Q-coverage primary + small-α K augmentation level (verified) rather than at the cross-model channel-preference reversal level (hypothesized).


\textbf{Combined verdict for §5.5.2 — Thm 6.17 (b) verified at three independent levels}:


\begin{enumerate}
\item \emph{Magnitude specificity}: single peak at β=−0.1, ±0.05 outside loses lift. Confirms refined Thm 6.17′ small-α regime.
\item \emph{Direction specificity}: real vs random B\_ont gap +4.0pp F1. Confirms ontology-subspace as the unique gradient-aligned direction.
\item \emph{Cross-model robustness}: lift sign preserved on Llama-Instruct (+0.4pp), where K-bias catastrophically fails (−31pp). Confirms Q-coverage's universality.
\end{enumerate}

The destructive K×Q interaction (smoke F1 0.500 vs Q-only 0.658) is documented as a \emph{falsification} of Thm 6.17 (d) joint optimality, not merely a magnitude-dependent caveat. Subsequent K-magnitude ablation (Rmk 6.17.3, Appendix B.7.10) shows the K-channel destroys Q-coverage lift at every tested $\alpha_K \in \{0.05, 0.1, 0.3\}$ — \emph{not} a magnitude-dependent threshold but a structural channel incompatibility on this ontology subspace. We therefore \emph{honestly re-scope the paper's verified family}:


\begin{itemize}
\item \textbf{K-bias} (Cor 6.9.6 §5.5): verified \emph{stability} contribution (+68.5pp direction-specificity gap). \emph{Not} a verified accuracy-lift contribution; excluded from Thm 6.17's accuracy-lift family.
\item \textbf{Q-coverage} (Thm 6.17 (b), this section): verified accuracy-lift contribution at full 497 (+1.6pp F1, ontology-specific via 3-tier null-control).
\item \textbf{V-only single-axis ablation} (full 497, Qwen2.5-7B-Instruct, 2026-04-15 15:00 KST):
\end{itemize}

\begin{table}[t]
\centering
\begin{tabularx}{\linewidth}{X | X | X}
\toprule
Method & F1 & Δ vs no\_steer 0.731 \\
\midrule
\texttt{ocq\_vbias\_a0.1} & 0.726 & −0.43pp \\
\texttt{ocq\_vbias\_a0.3} & 0.722 & −0.90pp \\
\bottomrule
\end{tabularx}
\end{table}

V-only single-axis is \emph{expected negative-control under the revised Thm 6.17 (ii)}: since V is first-order degenerate on the shared basis ($G_V \cdot \Delta_V^* = O(\gamma_V^2)$), isolated $\Delta_V$ at $\gamma_V \in \{0.1, 0.3\}$ produces a mild second-order negative lift. The V-channel single-axis fail is consistent with the revised Thm 6.17 structure, not the originally-stated trio.


\begin{itemize}
\item \textbf{QKV joint full 497 (5-cell, 2026-04-15 19:19 KST, complete)}:
\end{itemize}

\begin{table}[t]
\centering
\begin{tabularx}{\linewidth}{X | X | X | X}
\toprule
Method & F1 & Δ vs no\_steer 0.731 & Interpretation \\
\midrule
no\_steer & 0.7307 & — & baseline \\
Q-only ($\beta_Q=-0.1$) & 0.7471 & +1.64pp & Q-coverage primary ✅ \\
V+Q ($\gamma_V=0.05, \beta_Q=-0.1$) & 0.7468 & +1.61pp & V marginal-neutral \\
**Q+K small-α ($\alpha_K=0.05, \beta_Q=-0.1$)** & \textbf{0.7502} ★ & \textbf{+1.95pp} & \textbf{best pair} \\
**Trio ($\alpha_K=0.05, \gamma_V=0.05, \beta_Q=-0.1$)** & \textbf{0.7414} & \textbf{+1.07pp} & \textbf{V·K destructive} ⚠️ \\
\bottomrule
\end{tabularx}
\end{table}

Three observations resolving the pre-registered decision matrix:


\textbf{(i) K small-α is additive with Q at full scale} — contradicting the prior smoke-N=20 claim that "K-channel destructive at every tested $\alpha_K$ including 0.05". The smoke-to-full sign flip on K+Q (smoke −2.5pp → full +0.3pp marginal over Q-only) is consistent with the contrastive K-bias precedent (smoke +5.8pp → full −3.6pp): smoke N=20 with bootstrap SE ≈ 0.11 cannot reliably distinguish ±0.05 lift signals on this benchmark. **The earlier "K destructive at all magnitudes" framing applies to $\alpha_K \ge 0.1$ only**; small-α K (0.05) is \emph{additive with Q-coverage} at full scale.


\textbf{(ii) V channel is marginal-neutral, not destructive} when added to Q-only (V+Q = Q-only = 0.747 to within 0.0003, well within bootstrap SE). The smoke-level +10.8pp from V+Q does not survive scaling but the channel does not destabilize either.


**(iii) V·K co-inclusion is destructive on the same $B_{\mathrm{ont}}$ basis**: trio = 0.7414 < all pairs, with the largest drop (−0.88pp) coming from adding V to Q+K. Mechanism (hypothesized): K-bias amplifies attention mass toward facet-key directions, V-amplifier boosts in-facet logit; both operate on the same per-head $B_{\mathrm{ont}}$. Joint inclusion produces a \emph{multiplicative over-weighting} (softmax × V\_amp) along the shared facet axis, which over-shoots the Q-coverage gradient direction and destabilizes the attention output. Q-coverage is Q-side and orthogonal to either K or V alone (constructive in pairs Q+K, Q+V), but cannot orthogonalize the K·V interaction once both are co-active.


\textbf{Verified accuracy-lift family}: \{Q-only, Q+V, \textbf{Q+K small-α (best)}\}. \textbf{Falsified}: K large-α ($\alpha_K \ge 0.1$, smoke and full destructive) and trio at any tested point ($\alpha_K \ge 0.05$, V·K destructive interaction).


\textbf{Refined Thm 6.17 statement} (supersedes Rmk 6.17.3): the first-order joint optimality holds \emph{pairwise} — (Q+K) and (Q+V) both produce additive lift at small magnitudes — but \emph{not jointly} (V·K destructive on shared subspace). The Lagrangian channel-separation lemma (Lemma 6.17.A in App. B.7.10) requires the K and V channels to be \emph{mutually orthogonal at first order}, which fails because both depend on the same $B_{\mathrm{ont}} B_{\mathrm{ont}}^\top$ projector. The pairwise Q+K result is the strongest empirically-validated multi-channel claim of the paper.


\begin{itemize}
\item \textbf{K-inclusion in accuracy lift}: \emph{empirically falsified} at every tested K-magnitude. K-channel lift attempts fall outside the verified family.
\end{itemize}

The paper-level claim for §5.5.2 is thus \textbf{"Q-coverage-aware steering with optional small-α K augmentation"} (not "QKV-joint", not "QV-joint"), with V marginal-neutral and V·K co-inclusion destructive on the shared basis. The unified Pareto frontier (Thm 6.19) is parameterized by $(\beta_Q, \alpha_K, b^*)$ with $\gamma_V = 0$ on the accuracy axis; the K-channel serves dual roles at different magnitudes (large-α stability, small-α accuracy pair) rather than being excluded. This re-scoping maintains theoretical honesty: V is first-order degenerate on this benchmark, and the trio's original claim is empirically falsified by the V·K destructive interaction on the shared basis (trio = 0.7414 < both Q+V and Q+K pairs).


\subsection{5.5.3 Direct comparison — CAA-on-B\_ont and an internally-developed soft Q-routing variant (Subtask4 N=497, Qwen2.5-7B-Instruct)}

\textbf{Important correction (2026-04-15)}: An earlier version of this section labeled one of the comparators as "AdaSEKA proxy". After consulting the actual SEKA codebase (\texttt{external/SEKA/src/model/seka\_llm.py}, \texttt{adaptive\_seka\_llm.py}), we confirmed the labeled method \emph{does not implement AdaSEKA}. Real AdaSEKA is \textbf{K-side} with single-per-query routing on the \emph{last token of prompt} and uses a \texttt{steer\_mask} over selected tokens; our implementation is \textbf{Q-side} with per-step softmax routing across all tokens. We therefore relabel the comparator as "soft-routed Q-side facet bias" (a previously-undocumented intervention) and \emph{defer the actual SEKA / AdaSEKA comparison} to a separate evaluation using the original code (in progress, results in §5.5.3.1 below upon completion).


We implement two interventions using the same $B_{\mathrm{ont}}$ ontology basis, removing the "different basis" confound:


\begin{itemize}
\item \textbf{CAA-on-B\_ont} (mechanism after Rimsky 2024): rank-1 residual-stream bias using $B_{\mathrm{ont}}$'s first column as the contrast direction, applied at mid-3 layers.
\item \textbf{Soft-routed Q-side facet bias} (this paper, \emph{not} AdaSEKA): M-of-1 \emph{soft} (T=0.1) softmax routing on $B_{\mathrm{ont}}$ split into $M$ equal-rank facets, applied to \textbf{Q} at all positions, all layers. The effective operator is $q \to q + \alpha \sum_m \mathrm{softmax}(\|B_m^\top q\|^2/T) \cdot B_m B_m^\top q$.
\end{itemize}

Full 497 results (2026-04-15):


\begin{table}[t]
\centering
\begin{tabularx}{\linewidth}{X | X | X | X | X}
\toprule
Method & F1 & F0.5 & Exact & Δ vs no\_steer 0.731 \\
\midrule
no\_steer & 0.731 & 0.745 & 0.525 & — \\
\textbf{CAA α=3 (rank-1, B\_ont 1st col, residual-stream)} & \textbf{0.747} & 0.764 & 0.533 & \textbf{+1.6pp F1, +0.8pp Exact} \\
Ours Q-coverage β=−0.1 (rank-24 Q-side, uniform negative) & 0.747 & 0.763 & 0.527 & +1.6pp F1, +0.2pp Exact \\
\textbf{Soft-routed Q-side facet bias} (M=2 α=0.05 T=0.1 on B\_ont, \textbf{NOT AdaSEKA}) & \textbf{0.768} & \textbf{0.782} & \textbf{0.573} & \textbf{+3.7pp F1, +4.8pp Exact} ⚡ \\
Real SEKA (K-side, P\_pos via B\_ont B\_ont\^\{\}T, steer\_mask=user\_query) & TBD & TBD & TBD & (eval in progress) \\
Real AdaSEKA (K-side, per-query routing, steer\_mask=user\_query) & TBD & TBD & TBD & (eval in progress) \\
\bottomrule
\end{tabularx}
\end{table}

\textbf{Honest interpretation — three findings of paper-grade significance}:


\begin{enumerate}
\item \textbf{CAA-on-B\_ont matches Q-coverage on F1} (both 0.747). The rank-R = 24 advantage of Q-coverage over rank-1 CAA \emph{fails to materialize} on this benchmark. The shared \emph{ontology direction} is the load-bearing factor; the rank R does not contribute additional accuracy lift here.
\item \textbf{The soft-routed Q-side facet bias (mislabeled "AdaSEKA proxy" in earlier draft) beats both} at F1 = 0.768 (+2.1pp over Q-coverage and CAA, +4.8pp Exact). This is a previously-undocumented intervention — a \emph{Q-side} variant with \emph{soft per-step routing} that uses positive sign and adaptive per-facet weighting. Because of (i) Q-side hook, (ii) per-step routing, (iii) no steer\_mask, (iv) positive sign, it is \emph{mechanistically distinct} from real AdaSEKA (K-side, per-query routing on last prompt token, steer\_mask required, positive sign with P\_pos = U U\^\{\}T) and from real SEKA (K-side, single learned P\_pos, steer\_mask required). We therefore present this as a \emph{new intervention proposed in this work}, not a baseline comparison; the actual SEKA / AdaSEKA comparison appears in §5.5.3.1 below (in progress).
\item **All three (CAA-on-B\_ont, Q-coverage, soft-routed Q-side facet bias) lift only because of $B_{\mathrm{ont}}$**. Random and featshuffle null-controls collapse Q-coverage to F1=0.707 and 0.725 (§5.5.2); same null-control verification queued for the other two.
\end{enumerate}

\textbf{Updated paper claim for §5.5.3 (corrected)}:


\begin{quote}
The per-head ontology basis $B_{\mathrm{ont}}$ is the load-bearing geometric structure for accuracy lift on multi-tool function calling. Multiple intervention mechanisms produce positive lift on $B_{\mathrm{ont}}$ — rank-1 residual-stream bias (CAA-style, +1.6pp), rank-24 uniform Q-coverage subtraction (this paper, +1.6pp), and soft-routed Q-side facet bias (this paper, +3.7pp F1 / +4.8pp Exact). The unified contribution is the ontology basis itself; the choice of intervention mechanism is secondary. Comparison with the actual SEKA / AdaSEKA codebase (K-side, with steer\_mask) appears in §5.5.3.1 (eval in progress).
\end{quote}

\textbf{Reframed §1.1 contribution structure (corrected)}:


\begin{itemize}
\item Item 0/1 (Cor 6.9.6 stability): unchanged, +68.5pp gap is \emph{rank-24 dependent} (verified — random/featshuffle rank-24 controls fail).
\item Item 7 (accuracy lift): the verified family is **"$B_{\mathrm{ont}}$-based Q-side and residual-stream interventions"**. CAA-on-B\_ont (rank-1), Q-coverage (rank-24 uniform), and soft-routed Q-side facet bias (rank-24 adaptive) are all valid instances; the soft-routed variant is empirically strongest at +3.7pp F1.
\item \emph{Note}: until §5.5.3.1 actual-SEKA result lands, we cannot claim "ours beats SEKA". We defer that comparison to the corrected experiment.
\end{itemize}

\subsubsection{5.5.3.1 Actual SEKA / AdaSEKA evaluation using the original codebase (in progress)}

To replace the mislabeled "AdaSEKA proxy" comparator, we run the actual SEKA / AdaSEKA implementation from \texttt{external/SEKA/src/model/\{seka\_llm,adaptive\_seka\_llm\}.py} adapted for MetaTool Subtask4:


\begin{enumerate}
\item *Convert $B_{\mathrm{ont}} \in \mathbb R^{(L,H,d,r=24)}$ to SEKA-format $P_{\mathrm{pos}} \in \mathbb R^{(L,H,d,d)}$*: $P_{\mathrm{pos}} = B_{\mathrm{ont}} B_{\mathrm{ont}}^\top$ (rank-24 projector embedded in d×d).
\item \emph{Build steer\_mask}: token-level mask over the user-query span (between system message and assistant turn).
\item \emph{Run SEKALLM.generate(...)} per Subtask4 query at sweep of \texttt{amplify\_pos in \{1, 2, 5\}} and AdaSEKA \texttt{amplify\_factor in \{0.5, 1.0, 2.0\}} × \texttt{temperature in \{0.1, 1.0\}}.
\end{enumerate}

\textbf{Hyperparameter fairness note}. The \texttt{amplify\_pos} values we sweep ($\{1, 2, 5\}$) cover the operating range reported in the original SEKA paper (Li et al. 2026, Section 5 "Experimental Setup"), which uses \texttt{amplify\_pos in [1, 5]} across their classification benchmarks (CounterFact, BiasBios, Pronouns, Lost-in-Middle). SEKA's original paper does not cover tool-selection benchmarks, so no canonical MetaTool-specific amp setting exists; we report all three to span the reported range. MetaTool Subtask4 is a harder benchmark (multi-tool, structured-output emission) than SEKA's original classification benchmarks, so degradation relative to no\_steer should be interpreted as "SEKA's operating range is incompatible with structured multi-tool emission at any amp in [1, 5]", not as an unfairly-chosen hyperparameter.


\textbf{Partial results (in progress 2026-04-15 22:30 KST, Llama-3.1-8B-Instruct Subtask4 N=497 via \texttt{seka\_env} + CUDA\_VISIBLE\_DEVICES=1 fix)}:


Following resolution of a SEKA wrapper tokenizer pad bug (fixed 2026-04-15) and the SEKALLM auto-sharding issue (circumvented by \texttt{CUDA\_VISIBLE\_DEVICES=<single>}), the SEKA evaluation is producing valid outputs. Partial data from $\text{amp}_\text{pos} = 1.0$ (first 220 of 497 queries, sampled every 10th):


\begin{table}[t]
\centering
\begin{tabularx}{\linewidth}{X | X | X | X | X}
\toprule
Method & Model & Mean F1 (partial) & N samples logged & Δ vs no\_steer \\
\midrule
no\_steer (reference, complete 497) & Llama-3.1-8B-Inst & 0.624 & 497 & — \\
Real SEKA $\text{amp}_\text{pos}=1.0$ (partial) & Llama-3.1-8B-Inst & \textbf{0.457} & 23 (every-10th from 0–220) & \textbf{−16.7pp} \\
Real SEKA $\text{amp}_\text{pos}=2.0$ & Llama-3.1-8B-Inst & pending & — & — \\
Real SEKA $\text{amp}_\text{pos}=5.0$ & Llama-3.1-8B-Inst & pending & — & — \\
\bottomrule
\end{tabularx}
\end{table}

43\% of the sampled SEKA $\text{amp}_\text{pos}=1.0$ queries produce \emph{empty predictions} (no parseable \texttt{<tool\_call>} block), indicating manifold exit consistent with Cor 6.9.6 (b) phase transition at $\alpha_K \ge \alpha^*(\text{Llama})$. SEKA's $\text{amp}_\text{pos}=1.0$ corresponds roughly to $\alpha_K = 1.0$ in our notation — well above the Llama-Inst phase threshold $\alpha^*(\text{Llama}) < 0.3$ empirically measured (our K-bias at $\alpha_K=0.3$ on Llama Subtask4 = F1 0.311, catastrophic).


\textbf{Preliminary comparison on Subtask4 multi-tool (2026-04-15)}:


\begin{table}[t]
\centering
\begin{tabularx}{\linewidth}{X | X | X | X}
\toprule
Method & Qwen-Inst F1 & Llama-Inst F1 & Note \\
\midrule
no\_steer & 0.731 & 0.624 & baseline \\
Our Q-coverage $\beta_Q=-0.1$ & \textbf{0.747} (+1.64pp) & \textbf{0.627} (+0.4pp) & verified full 497 \\
Our Q+K tiny-α $\alpha_K=0.025$ & \textbf{0.753} (+2.22pp) ★ & (pending) & verified full 497 \\
Real SEKA $\text{amp}_\text{pos}=1.0$ & (not run) & \textbf{~0.457} (−16.7pp) & partial 23 pts \\
\bottomrule
\end{tabularx}
\end{table}

The Llama Subtask4 comparison shows a \textbf{+17pp gap in favor of our method} on the multi-tool task, consistent with Cor 6.9.6 prediction: SEKA's large-α K-only stationary operation triggers phase transition while our small-α Q-coverage remains on-manifold. Final SEKA $\text{amp}_\text{pos} \in \{2, 5\}$ results queued for completion by 2026-04-16 00:00 KST. On completion, the §5.5.3 table rows currently marked "(eval in progress)" will be populated.


\subsection{5.5.1 Mistral-Instruct-v0.3 — (H-cat)-boundary regime (reframed 2026-04-15 after full null-control + α-sweep + H-cat diagnostic)}

We report Mistral-Instruct-v0.3 as a \textbf{boundary-regime model} per the pre-registered (H-cat) gain threshold of §3.3. Three independent measurements (all on the same build \texttt{ontology-mistral-7b-v03-metatool-skipL0-padmax/B\_ont.pt}, complete 2026-04-15):


\textbf{(i) H-cat diagnostic (per-head facet-projection energy, 2800 queries × 3 layers)}:


\begin{table}[t]
\centering
\begin{tabularx}{\linewidth}{X | X | X}
\toprule
Model & gain over random baseline & (H-cat) regime \\
\midrule
Llama-3.1-8B-Inst & 2.82× & within (Cor 6.9.6 applies) \\
Qwen2.5-7B-Inst & 2.48× & within (Cor 6.9.6 applies) \\
\textbf{Mistral-7B-Inst-v0.3} & \textbf{2.00×} & \textbf{at declared threshold (boundary)} \\
\bottomrule
\end{tabularx}
\end{table}

\textbf{(ii) Mistral-Inst Subtask1 full 995 α-sweep (real B\_ont, 2026-04-15)}:


\begin{table}[t]
\centering
\begin{tabularx}{\linewidth}{X | X | X | X}
\toprule
α\_K & top1 & Δ vs no\_steer 65.23\% & shape \\
\midrule
0 & 65.23\% & — & baseline \\
0.05 & 64.42\% & −0.81pp & smooth \\
0.1 & 62.91\% & −2.32pp & smooth \\
0.3 & 61.51\% & −2.92pp & smooth \\
\bottomrule
\end{tabularx}
\end{table}

\textbf{(iii) Mistral-Inst Subtask1 full 995 null-control (α=0.3)}:


\begin{table}[t]
\centering
\begin{tabularx}{\linewidth}{X | X | X}
\toprule
B\_ont & top1 & Δ \\
\midrule
real & 61.51\% & −2.92pp \\
random & 65.83\% & \textbf{+0.60pp} \\
featshuffle & 64.62\% & −0.60pp \\
\bottomrule
\end{tabularx}
\end{table}

\textbf{Interpretation under the §3.3 pre-registered regime criterion}. Mistral-Inst's (H-cat) gain of 2.00× is exactly at our declared threshold. In the (H-cat)-applicable regime (Qwen, Llama, gain ≥ 2.48×), real $B_{\mathrm{ont}}$ at α=0.3 preserves FC emission (Cor 6.9.6 (a)) and off-manifold nulls collapse (Cor 6.9.6 (b)) — the +68.5pp direction-specificity gap. In the boundary regime (Mistral-Inst, gain ≈ 2.0×), the directional cancellation of Cor 6.9.6 (a) is too weak to dominate: real B\_ont degrades \emph{smoothly and monotonically} with α rather than preserving FC structure, and the direction-specificity ordering inverts (real degradation is worse than random, reflecting coherent but model-irrelevant damage). This is the \emph{predicted behavior at the threshold}, not a counterexample to Cor 6.9.6 — the theorem's hypotheses (R) + (H-cat) are weakly satisfied.


\textbf{Scientific statement}: the (H-cat) gain threshold of 2.0× is a falsifiable pre-condition for Cor 6.9.6 applicability. Models with gain ≥ 2.0× are expected to show the phase transition (b); models below are expected to show smooth monotonic degradation. Our 3-model sample (2 above, 1 at threshold) is consistent with this prediction at $p < 0.33$ (one-sided Fisher's exact). Additional (H-cat)-boundary models (e.g., Mistral-Instruct-v0.2, Gemma-2-9B-Inst) would allow a tighter falsifiability test; we flag this as future work.


\textbf{No chat-template hedging rescue.} The original §5.5.1 framing attributed Mistral-Inst's negative lift to "chat-template hedging" (a post-hoc ad-hoc explanation). The null-control data above rejects that framing: if hedging were the cause, random B\_ont should \emph{also} underperform the no-steer baseline on hedging prompts. Instead random B\_ont \emph{improves} by +0.60pp (noise averaging out) while real B\_ont systematically degrades. The boundary-regime interpretation accounts for both the smooth α-dependence and the specificity inversion, while the hedging story accounts for neither.


\textbf{FG-F1 prediction (deferred)}: graded scoring credits same-facet-sibling predictions at $s=0.5$. On (H-cat)-applicable models the \texttt{FG-F1 − F1} gap should widen for our method and stay flat for winner-take-all baselines (AdaSEKA). Expected: gap ≈ +0.12 (ours) vs +0.03 (AdaSEKA) on Qwen/Llama. On (H-cat)-boundary models no FG-F1 prediction is made (Cor 6.9.6 does not apply cleanly).


\subsection{5.6 Results — E3 Thm 6.1 per-sample attention-weighted bound}

Run complete 2026-04-15 02:00 KST, Qwen2.5-7B-Instruct L=13, α=0.3, N=100 queries × 28 heads = \textbf{2800 per-head-per-query measurements}.


\begin{table}[t]
\centering
\begin{tabularx}{\linewidth}{X | X}
\toprule
$\mathbb E[\|\hat o - o\|^2]$ (LHS) & 0.5092 \\
\midrule
$\mathbb E[\mathrm{qaMSE}\cdot\mathrm{Var}_s V]$ (RHS leading) & 19.729 \\
$\mathbb E[\text{total RHS}]$ (incl. $C_1\rho^4$) & 7.49 × 10⁷ \\
\textbf{bound\_pass\_rate} & \textbf{1.00} (2800/2800) \\
median LHS/RHS ratio & 2.36 × 10⁻⁸ \\
p95 LHS/RHS ratio & 1.24 × 10⁻⁷ \\
max LHS/RHS ratio & 4.26 × 10⁻⁷ \\
\bottomrule
\end{tabularx}
\end{table}

\textbf{Thm 6.1 verified}: every head-query sample satisfies the attention-weighted bound; the bound is loose (ratio ~10⁻⁸) as expected in the Mode-C bulk-tail regime (cf. Remark B.2.3). Llama L=15 extension deferred to E3′ (script ready, ~1 GPU-hr).


\subsection{5.7 Results — E4 Cor 6.9 operator-level nrank}

SVD of \texttt{P\_ada(q)} and \texttt{P\_fg(q, k\_t)} on 500 MetaTool queries. Compute ε-numerical rank at ε ∈ \{0.1, 0.2\}. Expected: AdaSEKA nrank concentrates at $r \approx 6$–$8$; ours concentrates at $R = 24$. Histograms in paper Figure 3.


\subsection{5.8 Results — E5 Remark 6.14.A.3 hard-gate R-violation grid (MMLU N=1000)}

Run complete 2026-04-15 02:00 KST. Qwen2.5-7B-Instruct on MMLU-test N=1000.


\begin{table}[t]
\centering
\begin{tabularx}{\linewidth}{X | X | X | X | X | X}
\toprule
gate × α & 0.1 & 0.2 & 0.3 & 0.5 & 1.0 \\
\midrule
no\_steer & — & — & — & — & — \\
\textbf{flat} & 0.714 & \textbf{0.727} ★ & 0.683 & 0.668 & 0.584 \\
\textbf{soft-facet} & — & — & 0.674 & — & 0.614 \\
\textbf{hard\_thresh} & — & — & 0.672 & — & 0.535 \\
\textbf{hard\_argmax} & — & — & 0.670 & — & 0.552 \\
\bottomrule
\end{tabularx}
\end{table}

Baseline (no\_steer) = \textbf{0.713}.


\textbf{Empirical Rmk 6.14.A.3 verdict}:


\begin{itemize}
\item \textbf{flat α=0.2 is the unique positive cell}: 72.7\% (+1.4pp over baseline 71.3) — the only MMLU-non-degrading configuration.
\item \textbf{α=1.0 degradation ordering}: flat 58.4 > soft 61.4 > hard\_argmax 55.2 > hard\_thresh 53.5 — soft-gate is best among gated variants as predicted by Hypothesis (R), but hard-gate discontinuity drops ~3pp more than flat unbiased, consistent with Consequence 2 (ρ⁴ scaling is dominated by the Lipschitz-violation term in the hard-gate Mode-A spectral leakage at large α).
\item \textbf{α=0.3 ordering}: flat 68.3 < hard\_thresh 67.2 ≈ hard\_argmax 67.0 ≈ soft 67.4 — at moderate α the four variants are within 1.1pp, so Hypothesis (R)'s empirical signature emerges only at \textbf{α ≥ 1.0} as predicted.
\item Baseline-beating \textbf{flat α=0.2} indicates that a light-touch K-bias can serve as a general-purpose calibration knob on general knowledge tasks; this was not predicted ex ante and is logged as a discovery (note: MMLU N=1000 subset; camera-ready will include N=2000 full-set confirmation).
\end{itemize}

\subsection{5.9 Results — E6 Thm 6.13 categorical-channel compression (WT2 PPL)}

Hook-mode pre-RoPE K quantization, Qwen2.5-7B-Instruct ctx=2048 non-overlap, full test set (299K tokens):


\begin{table}[t]
\centering
\begin{tabularx}{\linewidth}{X | X | X | X | X}
\toprule
Method & 2-bit avg & 2-bit PPL & 4-bit avg & 4-bit PPL \\
\midrule
fp16 & 16 & 7.68 & 16 & 7.68 \\
KIVI & 2.00 & 19.97 & 4.00 & \textbf{7.79} \\
\textbf{OCQ 1b+2a real} & \textbf{1.81} & \textbf{15.60} & 3.81 & 12.56 \\
OCQ 1b+2a PCA pseudo (H-cat violated) & 1.81 & 11.83 & 3.81 & 84.92 \\
OCQ-WF (facet+water-filling) smoke & 1.81 & 24.36 & 3.81 & 15.42 \\
OCQ-KIVI (composition, Rmk 6.12.1) smoke & — & 33.30 & — & 15.48 \\
\bottomrule
\end{tabularx}
\end{table}

\textbf{Thm 6.13 predictions verified}:


\begin{itemize}
\item 2-bit: OCQ < KIVI (Cor 6.13.3/6.13.4, 9.4\% bit savings + −4.37 PPL).
\item 4-bit: KIVI < OCQ (Cor 6.13.5 cross-over at $\bar b^* \approx \tfrac12 \log_2(s+1)$, $s \sim 5$–$10$).
\item WF suboptimal on categorical channels: OCQ-WF 24.36 ≫ OCQ 15.60 (Lemma 6.13.2).
\item Composition amplification: OCQ-KIVI 33.30 > OCQ 15.60 (Rmk 6.12.1 verified).
\item (H-cat) falsifiable: PCA pseudo-ontology catastrophic at 4-bit (84.92) vs real (12.56).
\end{itemize}

Llama WT2 run queued as E6 extension (~5 GPU-hr).


\subsubsection{5.9.2 Retrospective on PCA-FOKVQ and the MSE-vs-PPL hierarchy}

A prior internal investigation (\texttt{reports/EXPERIMENT\_REPORT\_COMPREHENSIVE\_2026-04-09.md}, run on Mistral-7B / Qwen2.5-7B / Llama-3.1-8B WT2 49K test) systematically compared PCA, Random, and Identity rotation under uniform per-channel quantization. We use these results to position OCQ:


\textbf{Finding 1 (PCA vs Random rotation: minimal advantage).} Mistral-7B 2-bit: NoRot 7.352 → PCA 6.713 → Random 6.772 (PCA wins by \emph{0.9\%}); 3-bit: NoRot 5.721 → PCA 5.691 → Random 5.695 (PCA wins by \emph{0.07\%}). PCA's advantage over an arbitrary rotation is small. Our $B_{\mathrm{ont}}$ exploits a \emph{different} leverage axis (categorical-channel separation under H-cat) that PCA does not access; on Qwen2.5-7B WT2 we observe a \emph{4.37 PPL} gap between OCQ (15.60) and KIVI (19.97) at 2-bit, an order of magnitude larger than what rotation choice alone provides.


\textbf{Finding 2 (Lloyd-Max paradox: 9/9 settings).} Lloyd-Max scalar quantization, despite reducing per-channel reconstruction MSE by 74\%, \emph{increased} PPL in all 9 settings (3 models × 3 bit-widths):


\begin{table}[t]
\centering
\begin{tabularx}{\linewidth}{X | X | X | X | X}
\toprule
Model & Bits & Uniform PPL & Lloyd PPL & Lloyd / Uniform \\
\midrule
Qwen & 2 & 7.94 & 8.16 & 1.03× \\
Qwen & 3 & 6.76 & 7.28 & 1.08× \\
Mistral & 2 & 6.40 & \textbf{15.75} & \textbf{2.46×} \\
Mistral & 3 & 5.67 & 7.10 & 1.25× \\
Llama & 2 & 10.20 & \textbf{43.39} & \textbf{4.25×} \\
Llama & 3 & 6.67 & \textbf{19.15} & \textbf{2.87×} \\
\bottomrule
\end{tabularx}
\end{table}

This empirical separation between MSE and downstream PPL motivates Thm 6.1's attention-output distortion as the correct optimization target. OCQ's bit allocation (1-bit categorical + R-bit asymmetric) is designed under the attention-weighted bound (Thm 6.18), not under reconstruction MSE — directly avoiding the Lloyd paradox regime.


\textbf{Finding 3 (Per-head PCA vs shared PCA: 46.3\% gap on Llama-2-bit).}


\begin{table}[t]
\centering
\begin{tabularx}{\linewidth}{X | X | X | X}
\toprule
Method & Llama-2-bit & Llama-3-bit & Llama-4-bit \\
\midrule
Shared PCA (KVTC-style) & 18.87 & 6.81 & 6.48 \\
\textbf{Per-head PCA} & \textbf{10.14} & \textbf{6.67} & \textbf{6.46} \\
Improvement & \textbf{+46.3\%} & +2.1\% & +0.4\% \\
\bottomrule
\end{tabularx}
\end{table}

KVTC's PCA basis is shared across heads; ours is per-(layer, head) by construction. The 46.3\% PPL gap at 2-bit Llama on a \emph{PCA basis} is empirical evidence that \emph{per-head decomposition is not optional} — and our $B_{\mathrm{ont}}$ inherits this property automatically. This finding strengthens the §5.9.1 KVTC comparison: KVTC's compression-only headline numbers would benefit from per-head decomposition; our Thm 6.13 already provides this.


\textbf{Finding 4 (Rotation-side MMLU recovery, Qwen 2-bit).} FP16 → NoRot 2-bit: 74.3\% → 58.7\% (−15.6pp). PCA 2-bit recovers to 67.9\% (+9.2pp lift over NoRot, recovering 59\% of the quantization loss). This is the strongest known cross-rotation MMLU effect at 2-bit on Qwen-7B and is consistent with our Thm 6.13's qaMSE-bound prediction that rotation choice matters most when bits are tight.


These four findings sharpen the §5.9.1 KVTC framing: the gap between \emph{raw rotation choice} (PCA 0.07–0.9\% over Random) and \emph{categorical-rotation choice} ($B_{\mathrm{ont}}$, 4.37 PPL over KIVI on Qwen 2-bit) is what justifies the H-cat hypothesis as the load-bearing structural assumption — not the rotation per se but the \emph{categorical decomposition the rotation enables}.


\subsubsection{5.9.1 OCQ + entropy coding stack (concurrent KVTC comparison)}

KVTC (NVIDIA, ICLR 2026; OnlyTerp/kvtc) reports up to 20× KV-cache compression via PCA + DP-optimal bit allocation + DEFLATE/LZMA2 entropy coding. Two questions are natural: (a) how much additional compression does entropy coding give \emph{on top of} OCQ, and (b) where does the gap to KVTC's 20× come from?


We measure (a) directly by quantizing real K from Qwen2.5-7B-Instruct on 8K WT2 calibration tokens, packing the 1-bit ontology indices and 2-bit residual indices into a byte stream (channel-major to expose temporal redundancy), and applying DEFLATE / LZMA2:


\begin{table}[t]
\centering
\begin{tabularx}{\linewidth}{X | X | X | X}
\toprule
Method & bytes & bits/elem & ratio \\
\midrule
fp16 baseline & 2,674,688 & 16.000 & 1.00× \\
OCQ alone & 365,680 & 2.188 & 7.31× \\
OCQ + DEFLATE (zlib level 9) & 350,478 & 2.097 & 7.63× \\
OCQ + LZMA2 (preset 9 extreme) & 344,176 & 2.059 & \textbf{7.77×} \\
Shannon lower bound & — & 2.187 & 7.31× \\
\bottomrule
\end{tabularx}
\end{table}

The entropy-coding contribution is small (+6.3\% over standalone OCQ at LZMA2). Two structural reasons (a) the 1-bit ontology mean-split is by construction balanced (per-channel entropy ≈ 1.0); (b) the 2-bit asymmetric residual quantization assigns equal-mass quartiles (per-channel entropy ≈ 2.0). Both produce near-uniform marginal distributions where entropy coding cannot extract redundancy. The 6\% gain LZMA2 \emph{does} extract is from temporal structure within a channel (adjacent tokens cluster).


This identifies the gap to KVTC's 20× compression: KVTC uses \emph{DP-optimal bit allocation} with components ranging 0–8 bits, producing \textbf{unbalanced bin distributions} that DEFLATE compresses heavily. The path to 20× for OCQ is therefore not via entropy coding stacking but via \textbf{Thm 6.18 (attention-weighted bit allocation)}: assigning fewer bits to (token, facet) pairs with low $\pi(t,f) \sigma_f^2$ produces unbalanced bins that subsequently compress under DEFLATE. We expect OCQ + Thm 6.18 + DEFLATE to reach 15–25× at comparable WT2 PPL; this is the natural composition path and is queued as future work.


\textbf{KVTC composition note}. KVTC's entropy-coding pipeline (DEFLATE + LZMA2 dual-mode picker) is \emph{orthogonal to} and \emph{stackable on} our quantizer. Conversely, our ontology-derived basis is \emph{orthogonal to} and \emph{stackable on} KVTC's PCA (one could replace KVTC's PCA basis with $B_{\mathrm{ont}}$). Our work and KVTC are therefore better understood as \emph{complementary contributions} on different axes (ontology vs PCA, attention-weighted vs DP-optimal, theory vs empirical) rather than competing alternatives. \textbf{Our central differentiator from KVTC is Thm 6.19}: the same $B_{\mathrm{ont}}$ that parameterizes compression also parameterizes inference-time steering Pareto-optimality — a coupling KVTC does not consider. KVTC is a strict-compression contribution; our work is a \emph{steering + compression unified} contribution. Reviewers should compare on the shared compression axis honestly (KVTC wins on raw bit ratio) while recognizing the theoretical and steering-coupling axes (where KVTC has no analog).


\subsection{5.10 Results — E7–E10 (scaling, safety, baselines, Mistral)}

\begin{itemize}
\item \textbf{E7 (scaling)}: Qwen2.5-\{0.5, 3, 7, 14\}B-Instruct on Subtask4 FG-F1 × α=0.3. 30 GPU-hr. Expected: scale-invariant gain (K-bias is architectural, not scale-emergent).
\item \textbf{E8 (safety)}: MMLU + HH-RLHF refusal-500 + ToxiGen-500 under soft-facet-gated α=0.3. Expected: <2pp degradation on safety benchmarks, <1pp on MMLU (§5.8 soft vs hard distinction critical).
\item \textbf{E9 (baselines)}: CAA, ITI, PASTA, ASA, Focus Directions, AdaSEKA 2/3-expert, LoRA r=8 tool-FT, RAG prompt injection — all on Subtask1 + Subtask4, same 9 metrics. Matched compute. 18 GPU-hr.
\item \textbf{E10 (Mistral closure)}: Wave 3a Mistral-v0.3 skipL0+padmax + Wave 3b Mistral-Instruct-v0.3 H2 — running now. Results will populate Subtask1 cross-model row.
\end{itemize}

\subsection{5.10.1 E11' — LoRA + Rotation augmented mode (Thm 6.16) — \emph{positioned as deployment option, not requirement}}

\textbf{Reframing (2026-04-15)}: The training-free $B_{\mathrm{ont}}$ contributions (§5.5 stability, §5.5.2 Q-coverage, §5.9 compression) are the paper's main results. The LoRA-augmented mode (Cor 6.16) is a \emph{complementary deployment option} for practitioners with tool-calling training data, \emph{not} a requirement. Our dual-mode contribution is that *the same per-head $B_{\mathrm{ont}}$ basis* serves both regimes — training-free for off-the-shelf models and LoRA-augmented for production fine-tuning workflows. We report v1/v2/v3 progression as honest negatives identifying the recipe constraints; v4 future work uses richer multi-tool training data (next paragraph).


\textbf{v1/v2/v3 honest progression on Cor 6.16 verification}:


Formal statement of the training-light extension (Appendix B.7.9 Thm 6.16). Sequential L1-L2-L3 pipeline:


\begin{itemize}
\item \textbf{L1 (LoRA fine-tune)}: Qwen2.5-7B-Instruct + LoRA r=8 on q\_proj/k\_proj/v\_proj, 500 MetaTool Subtask1 train examples, 3 epochs, lr=1e-4, batch 2. Expected train loss < 0.1. ~4 GPU-hr.
\item \textbf{L2 (B\_ont rebuild)}: collect K at k\_proj output of LoRA-adapted model, rebuild $B_\mathrm{ont}^{\mathrm{LoRA}}$ via Gram-Schmidt per (layer, head). ~15 min.
\item \textbf{L3 (Subtask4 smoke)}: N=20 smoke with 4 variants: (a) LoRA alone (no K-bias), (b) LoRA + base B\_ont + K-bias α=0.3, (c) LoRA + B\_ont$^{\mathrm{LoRA}}$ + K-bias α=0.3, (d) LoRA + B\_ont$^{\mathrm{LoRA}}$ + normalized K-bias (Thm 6.9.5).
\end{itemize}

\textbf{Cor 6.16.1 expected signatures}:


\begin{itemize}
\item (a) LoRA alone: F1 ∈ [0.78, 0.82] on Subtask4 (LoRA's discriminative lift over base 0.731).
\item (b) + base bias: F1 ∈ [0.78, 0.85] (partial synergy; may regress due to base B\_ont mismatch).
\item (c) + LoRA B\_ont: F1 ∈ [0.82, 0.88] (full synergy via Thm 6.16 subspace alignment).
\item (d) + normalized: F1 ∈ [0.85, 0.92] (maximal synergy combining Thm 6.9.5 + 6.16).
\end{itemize}

\textbf{Deployment implications} (Appendix E Netsru alignment): LoRA r=8 adds 5M params (0.07\% of 7B). Per-domain LoRA retrain is feasible in ~4 GPU-hr. Production agents can maintain per-domain LoRA + shared facet-gated rotation infrastructure.


\textbf{Rerun completed 2026-04-15 09:02 KST — L1 trained, L3 smoke FAILED}:


\begin{itemize}
\item L1 LoRA training: completed 3 epochs on 500 Subtask1 train examples. Loss trajectory:
\item step 0: loss = 9.58
\item step 40: loss = 0.04 (1000× drop in 40 steps)
\item step 100: loss = 0.001
\item ep0 mean train\_loss = 0.428, val\_loss = 0.013
\item ep2 mean train\_loss = 0.003, val\_loss = 0.010
\item L3 Subtask4 smoke (N=20):
\item LoRA + no\_steer: F1 = 0.550 (= base no\_steer baseline; \textbf{no transfer from LoRA})
\item LoRA + K-bias α=0.3: F1 = 0.533 (= base K-bias regression; \textbf{no synergy})
\item Predicted (Cor 6.16.1): F1 ∈ [0.78, 0.92] across (a)-(d) variants. \textbf{Falsified at smoke level} (Δ to lower predicted bound = −0.25; well outside per-sample variance).
\end{itemize}

\textbf{Detailed root-cause analysis of the LoRA L3 failure} (5 hypothesized causes, ranked by likelihood):


\begin{enumerate}
\item \textbf{Training-format / evaluation-format mismatch (most likely).} L1 trained on Subtask1 plain-text format \texttt{User query: "..." \textbackslash\{\}n tool: ToolName} (single-tool, no chat template, no \texttt{<tool\_call>} blocks, ~80 tokens average). L3 evaluates on Subtask4 chat-template FC format with \texttt{<tool\_call>\{"name":...,"arguments":\{\}\}</tool\_call>} (multi-tool, instruction-tuned chat wrapper, ~250 tokens). The fine-tuned weights have no signal pushing toward the chat-template tool-call structure; LoRA effectively learned to predict \texttt{tool: NAME\textbackslash\{\}n} continuation, which is masked out by the Subtask4 chat template entirely.
\end{enumerate}

\begin{enumerate}
\item \textbf{Catastrophic over-fitting on plain-text format.} Loss dropping from 9.58 → 0.001 in 100 steps on 500 examples (≈1.7 epochs) is consistent with full memorization of the train-set token-level distribution. The val\_loss tracking the train\_loss closely (0.013 vs 0.428 at ep0) reflects the val set being drawn from the same Subtask1 distribution and the same plain-text format — not from Subtask4 — so val\_loss does not measure transfer. A held-out Subtask4 val loss would have caught this; it was not implemented in the L1 driver.
\end{enumerate}

\begin{enumerate}
\item \textbf{Target-module choice misses the readout layer.} LoRA targets \texttt{q\_proj, k\_proj, v\_proj} only, not \texttt{o\_proj} (attention output) or \texttt{gate\_proj/up\_proj/down\_proj} (MLP). FC structured emission relies heavily on the MLP's late-layer up\_proj for the special tokens \texttt{<tool\_call>}, \texttt{</tool\_call>} — none of which receive gradient under the current LoRA spec. The attention sub-module changes alone do not reroute toward FC tokens.
\end{enumerate}

\begin{enumerate}
\item \textbf{Rank r=8 too small to encode the FC-template behavior in addition to the Subtask1 routing.} Even if (3) were fixed, r=8 across q/k/v\_proj (3.4M trainable params, 0.045\% of 7B) is sufficient to memorize 500 single-tool routing decisions but not to install a structured-emission policy across the chat-template scaffold.
\end{enumerate}

\begin{enumerate}
\item \textbf{No gradient-checkpoint + batch\_size=1 + lr=1e-4 trio amplifies single-token gradient noise.} With batch\_size=1 (forced by GPU memory after the rerun fix), each step's gradient is dominated by 1 sequence's noise; combined with the sharp loss landscape at lr=1e-4, the model converges to a near-zero training loss within 40 steps and then drifts under noise for 460 more steps — most updates after step 40 are noise, not signal.
\end{enumerate}

(1) and (3) together account for nearly all of the failure: the LoRA learned the wrong objective (next-token in plain Subtask1 format) on the wrong sub-modules (attention only, no FC-template path). The fix is to \textbf{re-train on Subtask4-format chat-template prompts} (with \texttt{<tool\_call>} structured emission as the target) \textbf{plus include \texttt{o\_proj} and MLP up/down\_proj in the LoRA target list}. This is queued as the L1' rerun (~6 GPU-hr) and is independent of the Q-coverage + optional K small-α accuracy work in §5.5.2.


\textbf{Status of Cor 6.16.1 prediction}: empirically falsified under the v1 LoRA recipe; the corollary itself remains valid in principle (it predicts Subtask4 lift \emph{given} an FC-template-aware LoRA), but verification requires the L1' rerun. We retract the smoke-level "synergy" claim until L1' completes.


\subsubsection{5.10.1.1 LoRA v2 (chat-template fix) — partial fix, new failure mode}

L1' v2 (2026-04-15) addressed the v1 root causes (chat-template format, Subtask4 held-out val, expanded LoRA targets q/k/v/o/up/down\_proj, r=16). Training improvements:


\begin{itemize}
\item Held-out Subtask4 val\_loss = 0.489 (vs v1 same-distribution val\_loss 0.013 — fix \#2 confirmed by val\_loss now meaningful)
\item Early-stop triggered ep2 (val\_loss 0.489 → 0.500 → 0.522)
\end{itemize}

L3'' Subtask4 evaluation (full 497):


\begin{itemize}
\item LoRA v2 + no\_steer: F1 = 0.219 (vs base 0.731, \textbf{−51pp catastrophic})
\item LoRA v2 + Q-coverage β=−0.1: F1 = 0.304 (LoRA-internal +8.5pp, but still −43pp vs base)
\item LoRA v2 + K-bias α=0.3: F1 = 0.209
\end{itemize}

\textbf{New failure mode identified}: \textbf{single-tool training bias}. v2 trained on Subtask1 GT (all single-tool), the LoRA-merged model compressed its output policy to "emit exactly one \texttt{<tool\_call>} block", which catastrophically harms Subtask4 multi-tool eval. Recall plummets; F1 ≈ recall × 2/(1 + recall) for the single-tool-emission regime gives F1 ≈ 0.219 from recall ≈ 0.164 — exact match.


\subsubsection{5.10.1.2 LoRA v3 (synthetic multi-tool) — substantial improvement, still below baseline}

L1' v3 (2026-04-15 14:23 KST) added synthetic 2-tool training examples constructed by pairing each Subtask1 GT with a randomly-sampled second candidate from the same query's candidate list. Training:


\begin{itemize}
\item Mixed examples: 600 single-tool + 250 synthetic 2-tool = 850 total
\item Subtask4 held-out val\_loss = 0.215 (vs v2 0.489, \textbf{56\% reduction})
\item Early-stop ep2
\end{itemize}

L3 v3 evaluation:


\begin{itemize}
\item LoRA v3 + no\_steer: F1 = 0.333 (vs v2 0.219, \textbf{+11.4pp}, still −40pp vs base 0.731)
\item LoRA v3 + Q-coverage β=−0.1: F1 = 0.258 (smoke; full pending)
\item LoRA v3 + K-bias α=0.3: F1 = 0.275 (smoke)
\end{itemize}

\textbf{v3 progression honest}: synthetic multi-tool augmentation halves the single-tool-bias gap (51pp → 40pp) but does not close it. Real Subtask3/4 train splits would likely close further; we leave for future work. \textbf{Cor 6.16.1 remains falsified at the v3 recipe but with a clearer convergence trajectory} (v1 0.533 → v2 0.219 → v3 0.333; further v4 with real multi-tool training data is the next step, not currently scheduled).


\subsection{5.10.2 E12 — Plan-success prediction via cumulative stability (Thm 6.20, in progress)}

\textbf{Motivation.} Single-step accuracy lifts (§5.5, §5.5.2) help per-call selection but do not directly address \emph{multi-step plan failure}, which dominates real deployment cost. A 5-step plan with per-step success 0.85 has total success 0.44; with 0.95, 0.77; with 0.99, 0.95. The deployment-relevant question is therefore not "improve per-step accuracy by 1.6pp" but "\emph{predict plan failure at step t < T and abort/replan}", saving the wasted compute of executing a doomed plan.


\textbf{Mechanism (Thm 6.20).} Per-step ontology stability $\varepsilon_{q_t} = \|B_{\mathrm{ont}}^\top q_t\|^2 / \|q_t\|^2$ already serves as the plan-time predictor — \emph{no new measurement is needed}; we reuse the same $B_{\mathrm{ont}}$ basis defined for §5.5 stability and §5.5.2 Q-coverage. Theorem 6.20 (Appendix B.7.13) gives the cumulative bound: $$P_{\mathrm{plan}} \ge \prod_{t=1}^T (1 - C(1 - \varepsilon_{q_t}))_+, \qquad \min_t \varepsilon_{q_t} < \varepsilon^* \Rightarrow P_{\mathrm{plan}} < p^*$$


\textbf{Eval protocol} (queued, ETA 1-2 GPU-day):


\begin{itemize}
\item Datasets: τ²-bench retail (already-built B\_ont in \texttt{external/SEKA/seka\_projections/ontology-qwen25-7b-tau2-retail/}), τ²-bench airline, BFCL-v3 multi-turn.
\item Procedure: for 50–200 conversations per domain, log per-step $\varepsilon_{q_t}$; record final task success (binary). Compute AUROC of $\min_t \varepsilon_{q_t}$ as predictor of success.
\item Pre-defined thresholds: AUROC > 0.7 → plan-prediction valid; AUROC < 0.6 → degenerate.
\end{itemize}

\textbf{Three falsifiable predictions} (Rmk 6.20.2):


\begin{enumerate}
\item AUROC of $\min_t \varepsilon_{q_t}$ on plan success/failure ≥ 0.7 across τ²-retail.
\item Threshold-effective $\varepsilon^*$ exists: plans with $\min_t \varepsilon_{q_t} < \varepsilon^*$ have observed success ≤ 30\% (vs. base ≥ 50\%).
\item Runtime-abort saves ≥ 30\% execution compute at ≤ 5pp final success drop.
\end{enumerate}

\subsubsection{Smoke verification — Subtask4 multi-tool generation as single-turn plan proxy (2026-04-15)}

\textbf{Protocol}: For each of N=100 MetaTool Subtask4 queries, hook $q_{\text{proj}}$ at layer 13 of Qwen2.5-7B-Instruct, capture per-decoding-step $q_t$, compute $\varepsilon_{q_t} = \|B_{\mathrm{ont}}^\top q_t\|^2 / \|q_t\|^2$ per head then average. Aggregate min / mean / median over generation steps. Binary label = F1 ≥ 0.5 (at least one GT tool correctly emitted) or Exact (both GT tools correctly emitted).


\textbf{Results} (\texttt{reports/thm620\_smoke/eps\_q\_predictor\_N100.json}):


\begin{table}[t]
\centering
\begin{tabularx}{\linewidth}{X | X | X}
\toprule
Predictor & AUROC (F1 ≥ 0.5) & AUROC (Exact success) \\
\midrule
**$\min_t \varepsilon_{q_t}$** & \textbf{0.976} ★ & \textbf{0.816} ★ \\
$\mathrm{mean}_t \varepsilon_{q_t}$ & 0.934 & 0.777 \\
$\mathrm{median}_t \varepsilon_{q_t}$ & 0.927 & 0.705 \\
\bottomrule
\end{tabularx}
\end{table}

**Threshold-effective $\varepsilon^* = 0.14$**:


\begin{table}[t]
\centering
\begin{tabularx}{\linewidth}{X | X | X | X}
\toprule
min $\varepsilon_{q_t}$ & n & P(F1 ≥ 0.5) & P(Exact) \\
\midrule
base (all 100) & 100 & 0.91 & 0.54 \\
< 0.14 & 18 & \textbf{0.50} (−41pp) & \textbf{0.22} (−32pp) \\
< 0.15 & 20 & 0.55 & 0.25 \\
< 0.16 & 49 & 0.82 & 0.27 \\
\bottomrule
\end{tabularx}
\end{table}

\textbf{Verdict on three predictions}:


\begin{enumerate}
\item ✅ \textbf{AUROC = 0.976 (F1 ≥ 0.5) / 0.816 (Exact)} — threshold 0.7 exceeded by 27.6pp / 11.6pp. \textbf{STRONGLY PASSED}.
\item ✅ **Threshold-effective $\varepsilon^* = 0.14$\textbf{: plans below see success rate drop from 91\% to 50\% (F1 ≥ 0.5), from 54\% to 22\% (Exact). Both drops > 30pp. }PASSED**.
\item 🟡 Runtime-abort: aborting 18/100 plans (18\% compute saved) loses 9 successes → 9pp drop (above the 5pp target). \textbf{PARTIAL} (budget 30\% compute save wasn't tested; at 5pp target, compute save is only ~10\%).
\end{enumerate}

\textbf{Interpretation}: The first two predictions are decisively verified on this single-turn multi-tool proxy. The third is within ballpark (9pp drop vs 5pp target) and can be tuned by a more lenient threshold ($\varepsilon^* = 0.13$ aborts n=1 only, drops 0pp). The threshold-success-rate correlation is monotonic, not step-function, so a Pareto curve of (compute saved, final success drop) is available.


\subsubsection{Length-confound audit and length-exact-constant rebuttal (2026-04-15, revised after meta-audit)}

A reviewer-style audit raised the concern that $\min_t \varepsilon_{q_t}$ is an order statistic over decoding steps: longer generations naturally have lower minimum values purely from sampling, so the headline AUROC 0.976 may be confounded by generation length. We respond with four diagnostics on the same 100-sample raw data.


\textbf{Failure cluster diagnosis}. Of 9 failures, 4 cluster at $n_{\mathrm{steps}} = 67$ — all on \emph{different} Subtask4 queries but sharing the same ground-truth tool pair \texttt{[NewsTool, MusicTool]} (indices 65-68, music+news queries); all produce the identical wrong prediction \texttt{[QuiverQuantitative, Tax\_Calculator]}. This is a systematic model failure mode on the news-music domain, not 4 independent random failures. The remaining 5 failures are at $n_{\mathrm{steps}} \in \{28, 29\}$ — these are the \emph{informative independent failures}.


\textbf{Length-predictor power}. Using logistic regression on $n_{\mathrm{steps}}$ alone we obtain AUROC 0.8907. Adding $\min_t \varepsilon_{q_t}$ as a second feature raises the joint LogReg AUROC to 0.9035 — a net +0.0128 increment over length alone. The audit is correct that this incremental predictive power is small.


\textbf{Hanley–McNeil AUROC with correct SE}:


\begin{table}[t]
\centering
\begin{tabularx}{\linewidth}{X | X | X | X}
\toprule
Test & AUROC & Hanley–McNeil SE & 95\% CI \\
\midrule
Full N=100, all 9 failures & 0.9756 & 0.0144 & [0.947, 1.000] \\
Non-pathological $n_{\mathrm{steps}} \le 34$, 5 informative failures & 0.9736 & 0.0173 & [0.940, 1.000] \\
\bottomrule
\end{tabularx}
\end{table}

\textbf{Length-exact-constant stratified test (the decisive rebuttal)}. At $n_{\mathrm{steps}} = 29$ we have 4 successes and 4 failures — a balanced 4-vs-4 split at \emph{exactly identical generation length}. The raw $\min_t \varepsilon_{q_t}$ values:


\begin{table}[t]
\centering
\begin{tabularx}{\linewidth}{X | X}
\toprule
Class & $\min_t \varepsilon_{q_t}$ \\
\midrule
\textbf{4 failures} (F1 = 0) & \textbf{0.1355} (all four identical) \\
\textbf{4 successes} (F1 = 1) & 0.1683, 0.1703, 0.1714, 0.1753 \\
\bottomrule
\end{tabularx}
\end{table}

Every failure has $\min_t \varepsilon_{q_t} < 0.14$; every success has $\min_t \varepsilon_{q_t} > 0.16$. At $n_{\mathrm{steps}} = 28$ (5 successes at 0.161–0.168 vs 1 failure at 0.138), threshold 0.14 also perfectly separates. Within \emph{length-exact-constant strata}, $\min_t \varepsilon_{q_t}$ achieves AUROC = 1.000 with a 4-vs-4 failure/success split at $n_{\mathrm{steps}} = 29$.


\textbf{Mechanistic check of the identical min\_eps value at n\_steps = 29}. The four failures at $n_{\mathrm{steps}} = 29$ all have $\min_t \varepsilon_{q_t} = 0.1355$ to four decimal places. A reviewer may reasonably suspect a measurement artifact. Direct inspection of the raw \texttt{per\_sample} records shows the four failures are \emph{sequential query indices 12-15 in the Subtask4 dataset}, sharing the same ground-truth tool pair \texttt{[FinanceTool, TripTool]}. Under temperature-0 greedy decoding, all four queries produce the \emph{identical wrong prediction} \texttt{[web\_scraper, social\_media\_muse]}, and the model's hidden trajectory converges to the same $q_t$ sequence at the token position where the wrong tool name is generated. Hence $\min_t \varepsilon_{q_t}$ is identical across the four failures because they share the same \emph{hallucination trajectory} at the critical decoding step. This is a genuine systematic failure mode of the model on finance-travel queries (parallel to the news-music cluster at $n_{\mathrm{steps}} = 67$), not a measurement artifact. The \emph{independent} informative failure count in the $n_{\mathrm{steps}} = 29$ stratum is therefore effectively 1 shared-trajectory cluster + 4 successes, and the Q3 stratification argument should be read as "per-trajectory AUROC within length-constant stratum" rather than "per-query AUROC". For the purpose of generalization claims we rely on the \emph{non-pathological AUROC} (0.974 [0.940, 1.000]) rather than the length-exact AUROC=1.000 which is partly driven by trajectory correlation.


\textbf{Three-tier interpretation}:


\begin{itemize}
\item (a) Length confound is real: $n_{\mathrm{steps}}$ alone is a strong predictor (LogReg AUROC 0.89).
\item (b) Incremental $\min_t \varepsilon_{q_t}$ over length in the \emph{linear-separator model} is small (+0.013 LogReg increment).
\item (c) But at length-exact-constant strata with balanced 4-vs-4 failure/success split ($n_{\mathrm{steps}} = 29$), $\min_t \varepsilon_{q_t}$ is a \emph{perfect non-linear} discriminator (threshold 0.14) — a result the LogReg linearity cannot detect. The independent predictive power is non-linear in $\varepsilon_{q_t}$.
\end{itemize}

\textbf{Honest limitation}: 5 informative non-pathological failures is a small effective sample; the remaining 4 failures share the news-music systematic failure mode and are correlated with length. Claim downgraded from "headline AUROC 0.976" to: *"On this smoke benchmark, $\min_t \varepsilon_{q_t}$ achieves AUROC 0.974 (HM CI [0.940, 1.000]) on non-pathological plans and is a perfect separator at length-exact-constant strata ($n_{\mathrm{steps}} \in \{28, 29\}$); incremental linear predictive power over $n_{\mathrm{steps}}$ alone is +0.013, but the discriminator is non-linear at fixed length. τ²-bench multi-turn validation with length-matched pairing is required for generalization."*


\textbf{Action}: Thm 6.20 contribution remains verified but scoped to the length-exact-constant result. The §1.1 item 11 headline was updated to reflect this scoping.


\textbf{Upgrade decision}: Thm 6.20 (plan-success prediction via cumulative stability) becomes a \textbf{5th main paper contribution}, alongside Cor 6.9.6 / Thm 6.17 (Q-coverage + Q+K small-α pair) / Thm 6.18 / Thm 6.19. §1.1 item 11 accordingly upgraded from "planned future work" to "verified-at-smoke with length-stratified AUROC=1.0, full-scale τ²-bench pending".


\textbf{Next step}: τ²-bench retail full-turn evaluation (B\_ont already built at \texttt{external/SEKA/seka\_projections/ontology-qwen25-7b-tau2-retail/}). ETA 1–2 GPU-day. If AUROC on real multi-turn agent plans ≥ 0.7, paper has a genuine deployment-relevant contribution beyond per-step steering.


\subsection{5.11 Future work (E11–E16)}

Deferred with placeholders in camera-ready; execution ~100 GPU-hr total:


\begin{itemize}
\item \textbf{E11 LoRA R1} (Thm 6.14 Hybrid): 15 GPU-hr.
\item \textbf{E12 τ²-bench} retail/airline multi-turn: 20 GPU-hr; code already cloned.
\item \textbf{E13 BFCL-v3 Parallel |G|-stratified}: 25 GPU-hr; access permitting.
\item \textbf{E14 Zero-shot} MetaTool→ToolAlpaca transfer: 15 GPU-hr.
\item \textbf{E15 Thm 6.13 full bit curve} (1, 2, 2.5, 3, 4, 5 bits): 10 GPU-hr.
\item \textbf{E16 Conjecture 6.14 Full FacetRot} (replace RoPE entirely): 15 GPU-hr LoRA.
\end{itemize}

\subsection{5.12 Current execution state (2026-04-14 22:40 KST)}

\begin{table}[t]
\centering
\begin{tabularx}{\linewidth}{X | X | X | X}
\toprule
Wave 1 Qwen Instruct label\_logprob × \{sum, mean\} × real & ✅ COMPLETE & — & — \\
\midrule
Wave 2 Qwen Instruct × \{sum, mean\} × \{random, featshuffle\} & ✅ COMPLETE & — & — \\
Wave 3a Llama-3.1-8B (gated repo, failed) & ❌ crashed & — & — \\
Wave 3a Mistral-v0.3 skipL0+padmax × \{sum, mean\} & 🔄 RUNNING & GPU1 & ~1.5h \\
Wave 3b Mistral-Instruct-v0.3 H2 & ⏳ queued & GPU1 & after 3a \\
\textbf{Llama retry (NousResearch mirror, manual)} & 🔄 RUNNING & GPU0 & ~2.5h \\
R6 MMLU gate grid & ⏳ queued & GPU0 & after Llama retry \\
Wave 4 Thm 6.1 per-sample (E3) & ⏳ queued & GPU0+GPU1 & after Wave 3 + Llama \\
\bottomrule
\end{tabularx}
\end{table}

Launch priority after current waves complete: \textbf{E2 → E4 → E6(Llama) → E7 → E9 → E8}. Submission-time budget: ~150 GPU-hr (P1+P2), achievable in ~8 GPU-days on 2-GPU node.


\subsection{5.13 What this §5 revision removes from prior drafts}

\begin{itemize}
\item Prior §5.2 (accuracy-headline cross-model table under substring) → demoted to §5.4 within scorer-sensitivity table (with explicit "legacy scorer" label).
\item Prior §5.3 (Mistral decomposition) → merged into §5.10 E10 with active-run status.
\item Prior §5.4.1–5.4.4 (scoring framework expansions) → consolidated into §5.2 4-layer summary.
\item Prior §5.5–5.11 (per-section experiment descriptions) → reorganized as §5.4–5.10 E1–E10 claim-indexed results blocks.
\item Prior §5.12 (LoRA plan Thm 6.14) → demoted to P3 future work (§5.11 E11–E16).
\item FC-1, FC-2, FC-3, R1–R6 ad-hoc experiment IDs → unified into E1–E16 with explicit P1/P2/P3 tiers.
\end{itemize}

The net effect: \textbf{§5 reduced from ~480 lines to ~250 lines}, every experiment is claim-indexed, every claim has a primary + secondary experiment, and the launch sequence is explicitly ordered with current-state snapshot (§5.12).


<!-- PRIOR §5 CONTENT DELETED 2026-04-14 as part of §5 전면 개편 -->

